\documentclass[11pt]{article}
\usepackage[margin=1.1in]{geometry}
\usepackage{amsmath,amssymb,amsthm,mathtools}
\usepackage{enumitem}
\usepackage{hyperref}
\usepackage{microtype}
\usepackage{booktabs}
\usepackage{mathrsfs}
\hypersetup{colorlinks=true,linkcolor=blue,citecolor=blue,urlcolor=blue,hypertexnames=false,pdftitle={Sukhatme Fixed Points and the Terminal Poisson Process},pdfauthor={Christopher D. Long}}
\newtheorem{theorem}{Theorem}[section]
\newtheorem{proposition}[theorem]{Proposition}
\newtheorem{lemma}[theorem]{Lemma}
\newtheorem{corollary}[theorem]{Corollary}
\newtheorem{conjecture}[theorem]{Conjecture}
\theoremstyle{remark}
\newtheorem{remark}[theorem]{Remark}
\newtheorem{warning}[theorem]{Warning}
\newtheorem{definition}[theorem]{Definition}
\newcommand{\Pp}{\mathbb P}
\newcommand{\Ee}{\mathbb E}
\newcommand{\Rr}{\mathbb R}
\newcommand{\Nn}{\mathbb N}
\newcommand{\Zz}{\mathbb Z}
\newcommand{\one}{\mathbf 1}
\newcommand{\TV}{\mathrm{TV}}
\newcommand{\Pois}{\operatorname{Pois}}
\newcommand{\Exp}{\operatorname{Exp}}
\newcommand{\PPP}{\operatorname{PPP}}
\newcommand{\Var}{\operatorname{Var}}
\newcommand{\Cov}{\operatorname{Cov}}
\newcommand{\diag}{\operatorname{diag}}
\newcommand{\dd}{\,d}
\newcommand{\eps}{\varepsilon}
\newcommand{\calD}{\mathcal D}
\newcommand{\calC}{\mathcal C}
\newcommand{\calE}{\mathcal E}
\newcommand{\calG}{\mathcal G}
\title{Sukhatme Fixed Points and the Terminal Poisson Process}
\author{Christopher D. Long}
\date{Mechanical extraction --- proof audit pending}
\begin{document}\maketitle
\begin{abstract}
For the singular Luce profile with rates \(n,n-1,\ldots,1\), this paper proves that the expected number of fixed points is asymptotic to \(e^{-1}\log n\), establishes a local Poisson point process for fixed points on logarithmic terminal coordinates, and proves all fixed factorial-moment asymptotics. The proof develops an anisotropic terminal saddle, tilted lattice local estimates, product-scale localization, and a death-process domination for non-separated configurations.
\end{abstract}
\tableofcontents
\section{Status and provenance}
This source is a mechanical extraction from the frozen 139-page Luce manuscript at repository commit \texttt{6178a99a8f2ee7b6416edd026e7780b9b8fa9fa0}. The mathematical blocks between the archival sentinels below are copied byte-for-byte from the frozen source. Only the standalone preamble, abstract, section wrappers, provenance notice, and cited bibliography entries have been added.

The extraction preserves every theorem hypothesis, identity, limiting statement, topology, quantifier, and proof step in the selected theorem package. It does not audit, repair, strengthen, weaken, or modernize the mathematics. Its status remains a working mechanical extraction until a separate adversarial audit is completed.

\section{Exponential-race input}

% BEGIN ARCHIVED BLOCK: exponential-race-input
\subsection{The exponential race}
Let \(\theta_1,\ldots,\theta_n>0\), and let \(X_i\sim\Exp(\theta_i)\) be independent.  For distinct labels \(i_1,\ldots,i_n\),
\[
\Pp(X_{i_1}<\cdots<X_{i_n})
= \prod_{r=1}^n
 \frac{\theta_{i_r}}{\theta_{i_r}+\theta_{i_{r+1}}+\cdots+\theta_{i_n}}.
\]
This is the standard exponential race identity and follows from the memoryless property.


% END ARCHIVED BLOCK: exponential-race-input
% BEGIN ARCHIVED BLOCK: gumbel-threshold-input
\subsection{Gumbel priorities}\label{subsec:gumbel-priorities}
A standard Gumbel random variable has distribution
\[
        \Pp(G\le y)=\exp(-e^{-y}),
        \qquad y\in\Rr .
\]
If \(E_i\sim\Exp(1)\) are independent, then \(-\log E_i\) are independent standard Gumbels.  Since
\[
        X_i\stackrel{d}=E_i/\theta_i,
\]
the exponential race may equivalently be generated by the shifted priorities
\[
        Y_i=-\log X_i=\log\theta_i+G_i,
        \qquad 1\le i\le n,
\]
ordered in decreasing order.  Thus
\[
        \sigma(i)=1+\#\{j:Y_j>Y_i\}.
\]
For later reference we record the exact threshold formulas.  The density and upper-tail function of \(Y_i\) are
\[
        f_i(y)=\theta_i e^{-y}\exp(-\theta_i e^{-y}),
        \qquad
        \Pp(Y_i>y)=1-\exp(-\theta_i e^{-y}).
\]

\begin{lemma}[Gumbel threshold identities]\label{lem:gumbel-threshold-identities}
For a single specified label \(p\),
\[
        \Pp(\sigma(p)=q)
        =\int_{-\infty}^{\infty}
        f_p(y)\,
        \Pp\left(\sum_{j\ne p}\one_{\{Y_j>y\}}=q-1\right)\dd y .
\]
More generally, let \(p_1,\ldots,p_\ell\) be distinct labels and let
\(q_1<\cdots<q_\ell\) be distinct target positions.  Put
\(P=\{p_1,\ldots,p_\ell\}\) and
\[
        S_a(y)=\sum_{j\notin P}\one_{\{Y_j>y_a\}},
        \qquad 1\le a\le\ell .
\]
Then
\[
\begin{aligned}
&\Pp(\sigma(p_a)=q_a,\ 1\le a\le\ell) \\
&\quad=\int_{y_1>\cdots>y_\ell}
        \prod_{a=1}^\ell f_{p_a}(y_a)\,
        \Pp(S_a(y)=q_a-a,\ 1\le a\le\ell)
        \dd y_1\cdots\dd y_\ell .
\end{aligned}
\]
\end{lemma}

\begin{proof}
Condition on the specified priorities.  For one label, the event
\(\sigma(p)=q\) is exactly the event that precisely \(q-1\) nonspecified
priorities exceed the observed value of \(Y_p\).  This gives the first
identity.

For the multi-point identity, the ordering of the target positions forces the
specified priorities to satisfy \(Y_{p_1}>\cdots>Y_{p_\ell}\).  Conditional on
\(Y_{p_a}=y_a\) in this chamber, exactly \(a-1\) specified priorities are above
\(y_a\).  Hence \(\sigma(p_a)=q_a\) is equivalent to
\(S_a(y)=q_a-a\), simultaneously for \(1\le a\le\ell\).  Integrating against the
independent specified-priority densities gives the formula.
\end{proof}

% END ARCHIVED BLOCK: gumbel-threshold-input
% BEGIN ARCHIVED BLOCK: analytic-input
\subsection{Analytic input}
We use standard estimates for sums of independent Bernoulli variables.  They are stated explicitly to make clear what is invoked.

\begin{lemma}[Point probability and Bernstein bounds]\label{lem:bernstein-point}
Let \(Y=\sum_{j=1}^N \xi_j\), where the \(\xi_j\) are independent, not necessarily identically distributed, Bernoulli variables, and set \(\mu=\Ee Y\), \(v=\Var Y\).  There is an absolute constant \(C\) such that
\[
        \sup_m \Pp(Y=m)\le \frac{C}{\sqrt{1+v}}.
\]
Moreover, there is an absolute constant \(c>0\) such that, for every \(u\ge0\),
\[
        \Pp(|Y-\mu|\ge u)\le
        2\exp\left[-c\min\left(\frac{u^2}{v+1},u\right)\right].
\]
\end{lemma}

\begin{proof}
The point bound is the classical concentration-function estimate for sums of Bernoulli variables.  It follows from Fourier inversion and the standard inequality
\[
        |1-p+pe^{it}|
        \le \exp\left[-\frac{2}{\pi^2}p(1-p)\operatorname{dist}(t,2\pi\mathbb Z)^2\right],
\]
where \(\operatorname{dist}(t,2\pi\mathbb Z)\) denotes distance from \(t\) to the nearest multiple of \(2\pi\).  Bernstein's inequality gives the second bound for bounded independent summands.  See Petrov \cite[Chapter VII]{Petrov1975} for both estimates.
\end{proof}

\begin{lemma}[Uniform lattice local CLT with a tilted local bound]
\label{lem:llt}
Fix a dimension $d$ and a support bound $B<\infty$.  Let
\[
        Y_{n,m}=\sum_{j=1}^{N_n}\zeta_{n,m,j},
        \qquad m\in\mathcal M_n,
\]
be sums of independent $\mathbb Z^d$-valued random vectors satisfying
$\|\zeta_{n,m,j}\|\le B$.  Write
\[
        \mu_{n,m}=\Ee Y_{n,m},
        \qquad \Sigma_{n,m}=\Var(Y_{n,m}),
\]
and suppose that, uniformly in $m$,
\[
        c s_n I_d\le\Sigma_{n,m}\le C s_n I_d,
        \qquad s_n\longrightarrow\infty.
\]
For $\lambda\in\mathbb R^d$, let $\Pp_\lambda$ denote the product exponential
tilt
\[
 \Pp_\lambda(\zeta_{n,m,j}=x)
 =\frac{e^{\lambda\cdot x}\Pp(\zeta_{n,m,j}=x)}
 {\Ee e^{\lambda\cdot\zeta_{n,m,j}}}.
\]
Assume that there is $\lambda_0>0$ for which the following two properties hold
uniformly for $\|\lambda\|\le\lambda_0$:
\begin{enumerate}[label=(\roman*)]
\item the tilted covariance matrices satisfy
\[
        c_0s_nI_d\le\Var_\lambda(Y_{n,m})\le C_0s_nI_d;
\]
\item for every $\varepsilon>0$ there are $\kappa_\varepsilon>0$ and
$n_\varepsilon$ such that
\[
 \left|\prod_{j=1}^{N_n}
 \Ee_\lambda e^{iu\cdot\zeta_{n,m,j}}\right|
 \le e^{-\kappa_\varepsilon s_n}
\]
whenever $n\ge n_\varepsilon$ and
$u\in[-\pi,\pi]^d$ has distance at least $\varepsilon$ from
$2\pi\mathbb Z^d$.
\end{enumerate}
Then, uniformly in $m$ and $z\in\mathbb Z^d$,
\begin{align*}
 \Pp(Y_{n,m}=z)
 &=(2\pi)^{-d/2}(\det\Sigma_{n,m})^{-1/2}\\
 &\quad\times
 \exp\left[-\frac12(z-\mu_{n,m})^T
 \Sigma_{n,m}^{-1}(z-\mu_{n,m})\right]
 +o(s_n^{-d/2}).
\end{align*}
Moreover,
\[
        \sup_z\Pp(Y_{n,m}=z)\le C_1s_n^{-d/2}.
\]
For every fixed $K<\infty$ and $\gamma<1/2$, there are $c_1,C_2>0$ such
that, uniformly when
\[
        \|z-\mu_{n,m}\|\le K s_n^{1/2+\gamma},
\]
one has
\[
 \Pp(Y_{n,m}=z)
 \le C_2s_n^{-d/2}
 \exp\left[-c_1\frac{\|z-\mu_{n,m}\|^2}{s_n}\right].
\]
All constants depend only on the fixed dimension, the support bound, and the
uniform constants in the hypotheses.
\end{lemma}

\begin{proof}
We suppress $n,m$ and write $Y,\mu,\Sigma,s$.  Let
$X_j=\zeta_j-\Ee\zeta_j$ and
\[
        \Psi(u)=\Ee e^{iu\cdot(Y-\mu)}.
\]
Since $\|X_j\|\le2B$ and
$\sum_j\Ee\|X_j\|^2=\operatorname{tr}\Sigma=O(s)$,
\[
        \sum_j\Ee\|X_j\|^3=O(s).
\]
Consequently, uniformly for $\|u\|$ sufficiently small,
\[
        \log\Psi(u)=-\frac12u^T\Sigma u+O(s\|u\|^3).
        \tag{2.1}
\]
A symmetrization argument also gives, on a fixed neighborhood of the origin,
\[
        |\Psi(u)|\le e^{-cs\|u\|^2}.
        \tag{2.2}
\]
Indeed, if $X_j'$ is an independent copy of $X_j$, then
$1-|\Ee e^{iu\cdot X_j}|^2=
\Ee[1-\cos(u\cdot(X_j-X_j'))]$; for small $u$ the right-hand side is
bounded below by a constant multiple of $\Var(u\cdot X_j)$, and multiplication
over $j$ yields (2.2).

Fourier inversion on $[-\pi,\pi]^d$ now proves the local expansion.  Choose
$L_s\to\infty$ with $L_s^3/\sqrt{s}\to0$ and split the integral into
\[
 \|u\|\le L_s/\sqrt{s},\qquad
 L_s/\sqrt{s}<\|u\|\le\varepsilon,
 \qquad \|u\|>\varepsilon.
\]
On the first region, (2.1), followed by the change of variables
$v=\Sigma^{1/2}u$, gives the Gaussian Fourier integral with an
$o(s^{-d/2})$ error.  The estimate (2.2) makes the second region
$o(s^{-d/2})$, and tilted aperiodicity with $\lambda=0$ makes the third
region exponentially small.  The oscillatory factor
$e^{-iu\cdot(z-\mu)}$ has modulus one, so every error is uniform in $z$.
Taking absolute values in the same inversion formula proves the concentration
bound.  This is the usual fixed-dimensional Fourier proof of the lattice local
central limit theorem; compare Petrov \cite[Chapter~VII]{Petrov1975}.

It remains to prove the Gaussian local upper bound.  Let
$h=z-\mu$ and define the centered cumulant generating function
\[
        K(\lambda)=\log\Ee e^{\lambda\cdot(Y-\mu)}.
\]
For $\|\lambda\|\le\lambda_0$, its Hessian is
$\nabla^2K(\lambda)=\Var_\lambda(Y)$ and is therefore between
$c_0sI_d$ and $C_0sI_d$.  If $\|h\|=o(s)$, strong convexity shows that the
minimum of $K(\lambda)-h\cdot\lambda$ over the ball
$\|\lambda\|\le\lambda_0/2$ is attained in its interior for all large $n$.
Thus there is a unique $\lambda_h$ in that ball satisfying
\[
        \nabla K(\lambda_h)=h,
        \qquad \|\lambda_h\|\le C\|h\|/s.
\]
The associated Legendre transform
$I(h)=\lambda_h\cdot h-K(\lambda_h)$ satisfies
\[
        I(h)\ge c\|h\|^2/s.
\]
For example, the last inequality follows by testing the variational formula
$I(h)=\sup_\lambda\{\lambda\cdot h-K(\lambda)\}$ at a sufficiently small
multiple of $h/s$ and using the tilted covariance upper bound.

Under $\Pp_{\lambda_h}$, the mean of $Y$ is $z$.  Repeating the absolute-value
Fourier estimate above, now using the two tilted hypotheses, gives
\[
        \sup_w\Pp_{\lambda_h}(Y=w)\le Cs^{-d/2}.
\]
The change-of-measure identity therefore yields
\[
 \Pp(Y=z)=e^{-I(h)}\Pp_{\lambda_h}(Y=z)
 \le Cs^{-d/2}e^{-c\|h\|^2/s}.
\]
The stated displacement range is $o(s)$ because $\gamma<1/2$, completing the
proof.
\end{proof}



% END ARCHIVED BLOCK: analytic-input
% BEGIN ARCHIVED BLOCK: sukhatme-fixed-point-theory
\section{Sukhatme fixed points}

We now prove the singular Sukhatme fixed-point mean, a local terminal Poisson point-process limit, and fixed-order factorial moment asymptotics.  Throughout this section
\[
        \theta_i=n+1-i.
\]
We call this the Sukhatme profile because the same decreasing rates \(n,n-1,\ldots,1\) occur in the classical exponential-spacing representation associated with Sukhatme and R\'enyi \cite{Sukhatme1937,Renyi1953}.  No external result from that theory is used below; the citations record only the historical source of the rate pattern and its standard order-statistic interpretation.
It is convenient to use the terminal rate coordinate
\[
        a=n+1-i.
\]
Thus rate coordinate \(a\) is the actual exponential rate of the corresponding clock, and it corresponds to label \(i=n+1-a\).  We write
\[
        T_a\sim\Exp(a),\qquad 1\le a\le n,
\]
for these independent clocks.  This is a change of notation from the regular-profile section, where the rates were denoted \(\theta_{n,i}\).  The fixed-point indicator in rate coordinates is
\[
        I_{n,a}=\one\left\{
        \sum_{b\ne a}\one_{\{T_b>T_a\}}=a-1
        \right\},
\]
and the total number of fixed points is
\[
        F_n=\sum_{a=1}^n I_{n,a}.
\]
Thus \(F_n\) is the fixed-point count \(C_1(\sigma_n)\) specialized to the singular Sukhatme profile, written in terminal rate coordinates.

The scale in this section is terminal rather than bulk.  The relevant rates satisfy
\(a\to\infty\), \(a=o(n)\), and their clock times are of order \(1/a\).  A single rate \(a\) contributes asymptotically \(e^{-1}/a\).  Summing over a logarithmic range gives the mean \(e^{-1}\log n\).  Joint probabilities factor for rates separated on logarithmic scale, which gives a local Poisson point process on the coordinate \(\log a\).  The proof parallels the regular-profile argument, but with terminal rather than bulk geometry: a saddle computation and local central limit estimate again produce the local entry probability, while the integrated saddle is anisotropic and the residual domination comes from the terminal death process rather than the bulk endpoint and collision envelopes.  The proof has three ingredients: finite-tail saddle geometry, an anisotropic terminal local central limit estimate, and a death-process product bound for non-separated tuples.

\begin{lemma}[Death-process product bound]\label{lem:death-bound}
For every finite set \(A\subset[n]\),
\[
        \Pp(I_{n,a}=1\text{ for all }a\in A)
        \le \prod_{a\in A}\frac{2}{a+1}
        \le 2^{|A|}\prod_{a\in A}\frac1a .
\]
\end{lemma}

\begin{proof}
Expose the exponential race as a pure death process.  When the alive set is \(S\), the next deleted rate is \(b\in S\) with conditional probability
\[
        \frac{b}{\sum_{c\in S}c}.
\]
The event \(I_{n,a}=1\) says that rate \(a\) is deleted at the unique transition at which the number of alive particles drops from \(a\) to \(a-1\).  Equivalently, immediately before \(T_a\) there are \(a\) alive rates, including \(a\), and immediately after \(T_a\) there are \(a-1\).  Thus \(I_{n,a}=1\) is exactly the event that rate \(a\) is deleted at this death transition.  Condition on the entire history before this transition.  If rate \(a\) is already dead, the conditional probability is zero.  If it is alive, the alive set has size \(a\) and contains \(a\).  The smallest possible total rate of such a set is
\[
        1+2+\cdots+a=\frac{a(a+1)}2,
\]
so the conditional probability that rate \(a\) is next deleted is at most
\[
        \frac{a}{a(a+1)/2}=\frac{2}{a+1}.
\]
Order the elements of \(A\) decreasingly and condition successively on the histories before the corresponding death levels.  Multiplying the conditional bounds gives the result.
\end{proof}

\subsection{The one-point terminal saddle}

We first record the elementary finite-tail estimates used throughout the terminal analysis.  They are stated in relative form because this is the form needed below: when \(x=o(n)\), the omitted infinite tail is negligible relative to \(x\), but it need not be \(O(1)\).

\begin{lemma}[Finite-tail geometric estimates]\label{lem:terminal-geometric}
Let \(x=x_n\to\infty\) and \(x=o(n)\).  Uniformly for \(s\) in compact subsets of \((0,\infty)\),
\[
        \sum_{b=1}^n e^{-bs/x}=\frac{x}{s}(1+o(1)),
\]
\[
        \sum_{b=1}^n \frac{b}{x}e^{-bs/x}=\frac{x}{s^2}(1+o(1)),
        \qquad
        \sum_{b=1}^n e^{-2bs/x}=\frac{x}{2s}(1+o(1)).
\]
Also
\[
        \sum_{b=1}^n b^2 e^{-bs/x}=O(x^3),
\]
uniformly for the same range of \(s\).  The same estimates remain true if a bounded number of summands is omitted.
\end{lemma}

\begin{proof}
For example,
\[
        \sum_{b=1}^n e^{-bs/x}
        =\frac{e^{-s/x}(1-e^{-sn/x})}{1-e^{-s/x}}.
\]
Since \(s\) is bounded away from zero and infinity, \(e^{-sn/x}=o(1)\) uniformly and \((1-e^{-s/x})^{-1}=x/s+O(1)\).  This gives the first estimate in relative form.

For the weighted sum, put \(q=e^{-s/x}\).  Then
\[
        \sum_{b=1}^n bq^b
        =
        \frac{q\bigl(1-(n+1)q^n+nq^{n+1}\bigr)}{(1-q)^2}.
\]
Since \(n/x\to\infty\), the truncation terms are exponentially small on the \(n/x\) scale, while
\[
        \frac{q}{(1-q)^2}=\frac{x^2}{s^2}(1+o(1))
\]
uniformly on compact \(s\)-intervals.  Dividing by \(x\) gives the second estimate.  The third estimate is the first estimate with \(s\) replaced by \(2s\).  For the final bound, compare with the infinite tail:
\[
        \sum_{b=1}^n b^2e^{-bs/x}\le \sum_{b=1}^{\infty}b^2e^{-bs/x}=O(x^3),
\]
uniformly on compact \(s\)-intervals bounded away from zero.  Omitting a bounded number of summands changes the first, second, and third displayed sums by \(O(1)\), \(O(1)\), and \(O(1)\), respectively, and changes the final sum by at most \(O(x^2)\).  Indeed,
\[
        \sup_{b\ge1}\frac{b}{x}e^{-bs/x}=O(1),
        \qquad
        \sup_{b\ge1}b^2e^{-bs/x}=O(x^2),
\]
uniformly for \(s\) in compact subsets of \((0,\infty)\).  These errors are negligible relative to the corresponding main scales, or harmless for the displayed \(O(x^3)\) bound.
\end{proof}

\begin{lemma}[One-point terminal saddle geometry]\label{lem:suk-one-saddle}
Assume \(a\to\infty\) and \(a=o(n)\).  Define
\[
        \mu_{n,a}(s)=\sum_{b\ne a}e^{-bs/a},
        \qquad
        v_{n,a}(s)=\sum_{b\ne a}e^{-bs/a}(1-e^{-bs/a}).
\]
Let \(s_{n,a}^*\) be the unique solution of
\[
        \mu_{n,a}(s_{n,a}^*)=a-1.
\]
Then, uniformly in this range,
\[
        s_{n,a}^*\to1,
        \qquad
        -\mu_{n,a}'(s_{n,a}^*)\sim a,
        \qquad
        v_{n,a}(s_{n,a}^*)\sim \frac a2 .
\]
Moreover, for all sufficiently large \(n\), uniformly for \(s\in[1/2,2]\),
\[
        |\mu_{n,a}(s)-(a-1)|\ge c a |s-s_{n,a}^*|
\]
whenever \(|s-s_{n,a}^*|\le 1/4\), with an absolute constant \(c>0\).
\end{lemma}

\begin{proof}
The function \(\mu_{n,a}\) is continuous and strictly decreasing, tends to \(n-1\) as \(s\downarrow0\), and tends to zero as \(s\to\infty\).  Hence the saddle exists and is unique.  By Lemma \ref{lem:terminal-geometric}, uniformly for \(s\) in compact subsets of \((0,\infty)\),
\[
        \mu_{n,a}(s)=\frac{a}{s}(1+o(1)),
        \qquad
        -\mu_{n,a}'(s)=\sum_{b\ne a}\frac{b}{a}e^{-bs/a}=\frac{a}{s^2}(1+o(1)),
\]
and
\[
        v_{n,a}(s)
        =\sum_{b\ne a}e^{-bs/a}-\sum_{b\ne a}e^{-2bs/a}
        =\frac{a}{2s}(1+o(1)).
\]
These estimates also locate the saddle.  Fix \(\varepsilon\in(0,1)\).  Uniformly,
\[
        \mu_{n,a}(1-\varepsilon)
        =\frac{a}{1-\varepsilon}(1+o(1))>a-1,
\]
whereas
\[
        \mu_{n,a}(1+\varepsilon)
        =\frac{a}{1+\varepsilon}(1+o(1))<a-1
\]
for all sufficiently large \(n\).  Since \(\mu_{n,a}\) is strictly decreasing, this implies
\(1-\varepsilon<s_{n,a}^*<1+\varepsilon\), and hence \(s_{n,a}^*\to1\).  The displayed derivative and variance asymptotics follow by substitution.

Finally, on \([1/2,2]\) the derivative satisfies \(-\mu_{n,a}'(s)\ge c a\) for all large \(n\).  The mean-value theorem gives the stated linear separation from the saddle.
\end{proof}

\begin{lemma}[Sukhatme one-point estimate]\label{lem:suk-one}
Let \(p_{n,a}=\Pp(I_{n,a}=1)\).  Uniformly for \(a\to\infty\) and \(a=o(n)\),
\[
        p_{n,a}\sim \frac1{ea}.
\]
Moreover, for all \(1\le a\le n\),
\[
        p_{n,a}\le \frac{2}{a+1}.
\]
\end{lemma}

\begin{proof}
The upper bound is Lemma \ref{lem:death-bound} with \(A=\{a\}\).  It remains to prove the asymptotic.

Conditioning on \(T_a=t\) and then putting \(t=s/a\) gives
\[
        p_{n,a}=\int_0^\infty e^{-s}
        \Pp(Y_{n,a}(s)=a-1)\dd s,
\]
where
\[
        Y_{n,a}(s)=\sum_{b\ne a}\xi_b(s),
        \qquad \xi_b(s)\sim\operatorname{Bernoulli}(e^{-bs/a})
\]
are independent.  Let \(s_{n,a}^*\) be the finite saddle from Lemma \ref{lem:suk-one-saddle}.

In the saddle window write
\[
        s=s_{n,a}^*+\frac{u}{\sqrt a},
        \qquad |u|\le a^{1/10}.
\]
For $s\in[1/2,2]$, a positive proportion of the indices $b\asymp a$ have
Bernoulli parameters $e^{-bs/a}$ in a fixed compact subinterval of $(0,1)$.
Under every bounded scalar exponential tilt these parameters remain uniformly
nondegenerate.  The total tilted variance is comparable to $a$, and the product
characteristic function has exponential decay away from $2\pi\mathbb Z$.
Thus all tilted hypotheses of Lemma \ref{lem:llt} hold uniformly with scale
$s_n=a$.  We use its local estimate in two forms.  First, on each fixed
\(u\)-ball, Lemma \ref{lem:llt}, together with Lemma
\ref{lem:suk-one-saddle}, gives
\[
        \Pp(Y_{n,a}(s)=a-1)
        =\frac{1}{\sqrt{\pi a}}e^{-u^2}(1+o(1)).
\]
The exponent is obtained from
\[
        a-1-\mu_{n,a}(s)
        =-\mu_{n,a}'(s_{n,a}^*)(s-s_{n,a}^*)+o(\sqrt a)
        =u\sqrt a(1+o(1)),
\]
and \(v_{n,a}(s_{n,a}^*)\sim a/2\).  Second, the Gaussian upper bound in the same one-dimensional form of Lemma \ref{lem:llt} gives
\[
        \Pp(Y_{n,a}(s)=a-1)
        \le C a^{-1/2}e^{-c u^2},
        \qquad |u|\le a^{1/10}.
\]
This bound justifies extending the fixed \(u\)-ball to the whole growing saddle window.

The contribution away from the saddle is negligible.  If \(s\notin[1/2,2]\), the mean of \(Y_{n,a}(s)\) is separated from \(a-1\) by a positive multiple of \(a\), so Bernstein's inequality gives an exponentially small contribution.  If \(s\in[1/2,2]\) but \(|s-s_{n,a}^*|>a^{-2/5}\), Lemma \ref{lem:suk-one-saddle} gives
\[
        |\mu_{n,a}(s)-(a-1)|
        \ge c a |s-s_{n,a}^*|.
\]
Since \(v_{n,a}(s)=O(a)\) uniformly on \([1/2,2]\), Bernstein's inequality gives
\[
        \Pp(Y_{n,a}(s)=a-1)
        \le C\exp[-c a(s-s_{n,a}^*)^2].
\]
The integral of this bound over \(|s-s_{n,a}^*|>a^{-2/5}\) is \(o(a^{-1})\).  Also \(e^{-s}=e^{-s_{n,a}^*}(1+o(1))\) uniformly for \(|u|\le a^{1/10}\).  Therefore
\[
\begin{aligned}
        p_{n,a}
        &\sim \int_{-\infty}^{\infty}
        e^{-s_{n,a}^*}
        \frac{1}{\sqrt{\pi a}}e^{-u^2}\frac{\dd u}{\sqrt a}  \\
        &\sim \frac{e^{-1}}{a\sqrt\pi}\int_{-\infty}^{\infty}e^{-u^2}\dd u
        =\frac1{ea}.
\end{aligned}
\]
\end{proof}

\begin{theorem}[Sukhatme fixed-point mean]\label{thm:suk-mean}
For Sukhatme weights \(\theta_i=n+1-i\),
\[
        \Ee F_n\sim e^{-1}\log n.
\]
\end{theorem}

\begin{proof}
Choose integer cutoffs
\[
        L_n=\lceil\log n\rceil,
        \qquad
        U_n=\lfloor n/\log n\rfloor .
\]
On this range \(L_n\to\infty\) and \(U_n=o(n)\).  Hence the uniform estimate in Lemma \ref{lem:suk-one} gives
\[
        \sup_{L_n\le a\le U_n}
        \left|ea\,p_{n,a}-1\right|\to0.
\]
Therefore
\[
        \sum_{L_n\le a\le U_n}p_{n,a}
        =e^{-1}(1+o(1))
        \sum_{L_n\le a\le U_n}\frac1a
        \sim e^{-1}\log n.
\]
The lower range contributes
\[
        \sum_{a<L_n}p_{n,a}
        \le 2\sum_{a<L_n}\frac1{a+1}
        =O(\log\log n)=o(\log n),
\]
and the upper range contributes
\[
        \sum_{a>U_n}p_{n,a}
        \le 2\sum_{a>U_n}\frac1{a+1}
        =O(\log\log n)=o(\log n).
\]
Thus \(\Ee F_n=\sum_a p_{n,a}\sim e^{-1}\log n\).
\end{proof}

\subsection{Separated terminal joint estimates}

The joint estimate is anisotropic: the covariance scale in the \(j\)-th terminal coordinate is \(a_j\), not a common scalar scale.  We therefore use a terminal local CLT adapted to logarithmically separated rates.  The one-point saddle above differentiated with respect to the scaled variable \(s=at\); in the joint saddle below we differentiate with respect to the clock time \(t\), which accounts for the change from derivative scale \(a\) to derivative scale \(a^2\).

\begin{lemma}[Joint terminal saddle geometry]\label{lem:suk-joint-saddle}
Fix \(r\ge1\) and \(\eta>0\).  Uniformly over integers
\[
        1\ll a_1<\cdots<a_r\ll n,
        \qquad a_{j+1}/a_j\ge1+\eta,
\]
where the ratio-separation condition is empty when \(r=1\).  Let \(A=\{a_1,\ldots,a_r\}\) and define
\[
        \mu_j(t)=\sum_{b\notin A}e^{-bt},
        \qquad 1\le j\le r.
\]
There is a unique \(t_j^*>0\) satisfying
\[
        \mu_j(t_j^*)=a_j-j.
\]
Moreover
\[
        a_jt_j^*\to1,
        \qquad
        -\mu_j'(t_j^*)\sim a_j^2,
        \qquad
        \mu_j''(t_j^*)=O(a_j^3),
\]
uniformly in \(j\).  Finally, for all large \(n\),
\[
        t_1^*>t_2^*>\cdots>t_r^*,
\]
and the ratios \(t_j^*/t_{j+1}^*\) are bounded below by a constant greater than one depending only on \(\eta\).
\end{lemma}

\begin{proof}
Fix \(j\).  For \(t=s/a_j\), Lemma \ref{lem:terminal-geometric} gives, uniformly for \(s\) in compact subsets of \((0,\infty)\),
\[
        \mu_j(s/a_j)=\frac{a_j}{s}(1+o(1)),
        \qquad
        -\mu_j'(s/a_j)=\frac{a_j^2}{s^2}(1+o(1)).
\]
The omitted set \(A\) has fixed size \(r\).  It changes the first sum by \(O_r(1)\).  For the derivative sum, each omitted term satisfies
\[
        b e^{-bs/a_j}\le C a_j
\]
uniformly for \(s\) in compact subsets of \((0,\infty)\), and hence the omitted terms change the derivative sum by \(O_r(a_j)=o(a_j^2)\).  Thus both displayed estimates remain valid after removing \(A\).

The equation \(\mu_j(t)=a_j-j\) has a unique solution, since \(\mu_j\) is continuous and strictly decreasing from \(n-r\) to \(0\), while \(0<a_j-j<n-r\) for all large \(n\).  To locate it, fix \(\varepsilon\in(0,1)\).  The first displayed estimate gives
\[
        \mu_j((1-\varepsilon)/a_j)>a_j-j,
        \qquad
        \mu_j((1+\varepsilon)/a_j)<a_j-j
\]
for all sufficiently large \(n\), uniformly in \(j\).  Therefore
\[
        1-\varepsilon<a_jt_j^*<1+\varepsilon,
\]
and \(a_jt_j^*\to1\).  Substituting \(t=t_j^*\) into the preceding derivative estimate, equivalently taking \(s=a_jt_j^*\to1\), gives \(-\mu_j'(t_j^*)\sim a_j^2\) uniformly in \(j\).  The second-derivative bound follows from
\[
        \mu_j''(t)=\sum_{b\notin A}b^2e^{-bt}=O(a_j^3)
        \qquad(t\asymp a_j^{-1}),
\]
again by the same finite-geometric estimate.  Since \(t_j^*\sim1/a_j\) and \(a_{j+1}/a_j\ge1+\eta\), the saddle times are eventually strictly ordered and uniformly separated multiplicatively.
\end{proof}

\paragraph{Terminal interval notation.}
For the remainder of this subsection, fix an admissible tuple
\[
        1\ll a_1<\cdots<a_r\ll n,
        \qquad a_{j+1}/a_j\ge1+\eta,
\]
put \(A=\{a_1,\ldots,a_r\}\), and consider ordered times
\(t_1>\cdots>t_r>0\).  Define
\[
        Z_j(t)=\sum_{b\notin A}\one_{\{T_b>t_j\}},
        \qquad z_j=a_j-j,
        \qquad \mu_j(t_j)=\Ee Z_j(t).
\]
Let
\[
        I_1=(t_1,\infty),\qquad
        I_j=(t_j,t_{j-1}]\quad(2\le j\le r),\qquad
        I_{r+1}=(0,t_r],
\]
and write
\[
        Y_1(t)=Z_1(t),\qquad
        Y_j(t)=Z_j(t)-Z_{j-1}(t)\quad(2\le j\le r).
\]
The corresponding interval-count target is
\[
        y_1=a_1-1,
        \qquad
        y_j=a_j-a_{j-1}-1\quad(2\le j\le r).
\]
Finally, put
\[
\begin{aligned}
        m_j(t)&=\Ee Y_j(t),\\
        H_j(t)&=z_j-\mu_j(t_j),\qquad H_0(t)=0,\\
        d_j(t)&=y_j-m_j(t)=H_j(t)-H_{j-1}(t),\\
        s_j&=a_jt_j,\qquad s_j^*=a_jt_j^*.
\end{aligned}
\]
The transformation between \(Z\) and \(Y\) is triangular with determinant
one, so \(\Pp(Z(t)=z)=\Pp(Y(t)=y)\).

\begin{lemma}[Coordinatewise Fourier concentration]\label{lem:coord-fourier}
Let
\[
        W=\sum_{\nu\in\mathcal I}W_\nu
\]
be a sum of independent \(\mathbb Z^r\)-valued random vectors, each supported in a fixed cube \([-M_0,M_0]^r\).  Let
\(A_1,\ldots,A_r\ge1\).  Suppose that, for every coordinate \(j\), there is a set
\(\mathcal I_j\subset\mathcal I\) with
\(|\mathcal I_j|\ge c_0A_j\), and that the sets
\(\mathcal I_1,\ldots,\mathcal I_r\) are disjoint.  Suppose also that, for every
\(\nu\in\mathcal I_j\), the law of \(W_\nu\) gives mass at least \(c_0\) to two atoms whose difference is \(e_j\), the \(j\)-th coordinate vector.  Then there are constants
\(c,C>0\), depending only on \(r,M_0,c_0\), such that the characteristic function
\(\phi_W(u)=\Ee e^{iu\cdot W}\) satisfies
\[
        |\phi_W(u)|
        \le \exp\left[-c\sum_{j=1}^r
        A_j\operatorname{dist}(u_j,2\pi\mathbb Z)^2\right],
        \qquad u\in[-\pi,\pi]^r.
\]
Consequently,
\[
        \sup_z\Pp(W=z)\le C\prod_{j=1}^r A_j^{-1/2}.
\]
The same conclusions hold uniformly after any exponential tilt
\[
        \Pp_\lambda(W_\nu=x)
        =\frac{e^{\lambda\cdot x}\Pp(W_\nu=x)}
        {\Ee e^{\lambda\cdot W_\nu}},
        \qquad \|\lambda\|_\infty\le \Lambda,
\]
with constants depending additionally on \(\Lambda\).
\end{lemma}

\begin{proof}
For a single summand \(Y\), write \(p_x=\Pp(Y=x)\).  If two atoms
\(x\) and \(x+e_j\) both have mass at least \(c_0\), then
\[
\begin{aligned}
        1-|\Ee e^{iu\cdot Y}|^2
        &=\frac12\sum_{y,z}p_yp_z
        |e^{iu\cdot y}-e^{iu\cdot z}|^2  \\
        &\ge c_0^2 |1-e^{iu_j}|^2
        \ge c\,\operatorname{dist}(u_j,2\pi\mathbb Z)^2 .
\end{aligned}
\]
After decreasing \(c\), this gives
\[
        |\Ee e^{iu\cdot Y}|
        \le \exp\{-c\operatorname{dist}(u_j,2\pi\mathbb Z)^2\}.
\]
Multiplying this estimate over the disjoint blocks \(\mathcal I_j\) gives the displayed characteristic-function bound.  Fourier inversion then gives
\[
\begin{aligned}
        \Pp(W=z)
        &\le (2\pi)^{-r}\int_{[-\pi,\pi]^r}
        \exp\left[-c\sum_{j=1}^r
        A_j\operatorname{dist}(u_j,2\pi\mathbb Z)^2\right]\dd u \\
        &\le C\prod_{j=1}^r A_j^{-1/2}.
\end{aligned}
\]
Under a tilt with \(\|\lambda\|_\infty\le\Lambda\), all probabilities of atoms in the fixed support are multiplied by factors bounded between positive constants depending only on \(M_0,r,\Lambda\).  Hence the two distinguished atom probabilities remain bounded below, with a possibly smaller constant, and the same proof applies.
\end{proof}

\begin{lemma}[Weighted chain comparison]\label{lem:weighted-chain-comparison}
Fix $r\ge1$ and $\eta>0$.  Let
$0<A_1<\cdots<A_r$ satisfy $A_{j+1}/A_j\ge1+\eta$, and put
$u_{r+1}=0$.  Then
\[
        \sum_{j=1}^r A_j u_j^2
        \le C_{r,\eta}\sum_{j=1}^r A_j(u_j-u_{j+1})^2 .
\]
Equivalently, if $d_j=H_j-H_{j-1}$ with $H_0=0$, then
\[
        \sum_{j=1}^r \frac{d_j^2}{A_j}
        \asymp_{r,\eta}
        \sum_{j=1}^r \frac{H_j^2}{A_j} .
\]
\end{lemma}

\begin{proof}
Write $\Delta_j=u_j-u_{j+1}$.  Then $u_j=\sum_{k=j}^r\Delta_k$.  By Cauchy's inequality with the geometrically decaying weights $A_j/A_k\le(1+\eta)^{-(k-j)}$,
\[
        A_j u_j^2
        \le C_{r,\eta}\sum_{k=j}^r A_k\Delta_k^2 .
\]
Summing over $j$ gives the first display.  The second follows by applying the first inequality to the inverse triangular map and its reverse; since $r$ is fixed and the ratios are bounded below, both triangular maps have uniformly bounded operator norms in the weighted norms.
\end{proof}

\begin{lemma}[Truncated triangular comparison]\label{lem:truncated-triangular-comparison}
Fix $r\ge1$, $\eta>0$, and $M<\infty$.  Let
$0<A_1<\cdots<A_r$ satisfy $A_{j+1}/A_j\ge1+\eta$.  If
$d_j=H_j-H_{j-1}$, $H_0=0$, and $|d_j|\le M A_j$ for every $j$, then
\[
        \sum_{j=1}^r
        \min\left\{\frac{d_j^2}{A_j},A_j\right\}
        \ge c_{r,\eta,M}
        \sum_{j=1}^r
        \min\left\{\frac{H_j^2}{A_j},A_j\right\}.
\]
\end{lemma}

\begin{proof}
Write $x_j=d_j/A_j$.  Then $|x_j|\le M$ and
\[
        \frac{H_j}{A_j}=\sum_{k\le j}\frac{A_k}{A_j}x_k,
        \qquad
        \frac{A_k}{A_j}\le (1+\eta)^{-(j-k)} .
\]
Since $\min\{y^2,1\}\le y^2$, Cauchy's inequality and the geometric
bound on the ratios give, for each fixed $j$,
\[
\begin{aligned}
        A_j\min\left\{\frac{H_j^2}{A_j^2},1\right\}
        &\le A_j\left(\sum_{k\le j}\frac{A_k}{A_j}x_k\right)^2  \\
        &\le C_{r,\eta} A_j\sum_{k\le j}\frac{A_k}{A_j}x_k^2
         = C_{r,\eta}\sum_{k\le j} A_k x_k^2 .
\end{aligned}
\]
Because $|x_k|\le M$, we have
$x_k^2\le C_M\min\{x_k^2,1\}$.  Summing the preceding inequality over
$j$ and using that $r$ is fixed gives
\[
        \sum_{j=1}^r A_j
        \min\left\{\frac{H_j^2}{A_j^2},1\right\}
        \le C_{r,\eta,M}
        \sum_{k=1}^r A_k\min\{x_k^2,1\}.
\]
This is exactly the claimed inequality after rewriting
$A_j\min\{H_j^2/A_j^2,1\}=\min\{H_j^2/A_j,A_j\}$ and
$A_k\min\{x_k^2,1\}=\min\{d_k^2/A_k,A_k\}$.
\end{proof}

\begin{lemma}[Boundary-aware truncated chain comparison]
\label{lem:boundary-truncated-chain}
Fix \(r\ge1\) and \(M<\infty\).  Let
\(0<A_1<\cdots<A_r\), let \(B_-\) and \(B_+\) be disjoint subsets of
\(\{1,\ldots,r\}\), and put
\[
        G=\{1,\ldots,r\}\setminus(B_-\cup B_+).
\]
Let \(H_0=0\), \(d_j=H_j-H_{j-1}\), and
\[
        \phi_j(x)=\min\left\{\frac{x^2}{A_j},A_j\right\}.
\]
Assume
\[
        |d_j|\le M A_j\quad(j\in G),
        \qquad
        |H_j|\le M A_j\quad(j\in B_+).
\]
Then
\[
\begin{aligned}
&\sum_{j\in G}\phi_j(d_j)
 +\sum_{j\in B_-}|H_j|
 +\sum_{j\in B_+}A_j\\
&\qquad\ge c_{r,M}
 \sum_{j\in G}\phi_j(H_j).
\end{aligned}
\]
\end{lemma}

\begin{proof}
Decompose \(G\) into maximal consecutive runs.  Let
\(\{\ell,\ell+1,\ldots,u\}\) be one such run and put
\(b=\ell-1\).  For every \(j\) in the run,
\[
        H_j=H_b+\sum_{k=\ell}^j d_k,
\]
where \(H_b=H_0=0\) if \(\ell=1\).  For a fixed positive number \(A\) and
at most \(r+1\) real numbers \(x_i\),
\[
        \min\left\{\frac{(\sum_i x_i)^2}{A},A\right\}
        \le C_r\sum_i
        \min\left\{\frac{x_i^2}{A},A\right\}.
        \tag{7.1}
\]
Indeed, if one term on the right equals \(A\), the claim is immediate; otherwise
it follows from Cauchy's inequality.

If \(k\le j\) belongs to the run, then \(A_k\le A_j\) and
\(|d_k|\le M A_k\).  Hence
\[
        \phi_j(d_k)\le C_M\phi_k(d_k).
        \tag{7.2}
\]
To see this, the assertion is immediate when \(|d_k|\le A_k\); when
\(|d_k|>A_k\), one has
\(\phi_k(d_k)=A_k\) and
\(\phi_j(d_k)\le d_k^2/A_j\le M^2A_k\).

It remains to control the left boundary.  If \(b\in B_-\), then
\[
        \phi_j(H_b)\le |H_b|.
\]
If \(b\in B_+\), then \(|H_b|\le M A_b\) and \(A_b\le A_j\), so
\[
        \phi_j(H_b)\le M^2 A_b.
\]
Combining these estimates with (7.1) and (7.2), and then summing first over
\(j\) in the run and finally over the finitely many runs, proves the lemma.
\end{proof}

\begin{lemma}[Edge Fourier concentration]\label{lem:edge-fourier}
Let $e_1,\ldots,e_r$ be the coordinate vectors in $\mathbb Z^r$ and put $e_{r+1}=0$.  Let
\[
        W=\sum_{\nu\in\mathcal I}W_\nu
\]
be a sum of independent random vectors supported on
$\{e_1,\ldots,e_r,0\}$.  Let $A_1,\ldots,A_r\ge1$ satisfy
$A_{j+1}/A_j\ge1+\eta$ for $j<r$.  Suppose that, for every edge $j$, there is a set
$\mathcal I_j\subset\mathcal I$ with $|\mathcal I_j|\ge c_0A_j$, the sets $\mathcal I_j$ are disjoint, and for every $\nu\in\mathcal I_j$ the two atoms $e_j$ and $e_{j+1}$ each have probability at least $c_0$.  Then
\[
        |\Ee e^{iu\cdot W}|
        \le \exp\left[-c\sum_{j=1}^r
        A_j\operatorname{dist}(u_j-u_{j+1},2\pi\mathbb Z)^2\right],
        \qquad u_{r+1}=0,
\]
and consequently
\[
        \sup_z\Pp(W=z)\le C\prod_{j=1}^r A_j^{-1/2}.
\]
The same conclusions hold uniformly after any exponential tilt
$\|\lambda\|_\infty\le\Lambda$, with constants depending additionally on
$\Lambda$.
\end{lemma}

\begin{proof}
If a summand $Y$ gives mass at least $c_0$ to $e_j$ and to $e_{j+1}$, then
\[
        1-|\Ee e^{iu\cdot Y}|^2
        \ge c_0^2|e^{iu_j}-e^{iu_{j+1}}|^2
        \ge c\operatorname{dist}(u_j-u_{j+1},2\pi\mathbb Z)^2 .
\]
As in Lemma \ref{lem:coord-fourier}, this implies
\[
        |\Ee e^{iu\cdot Y}|
        \le \exp\{-c\operatorname{dist}(u_j-u_{j+1},2\pi\mathbb Z)^2\}.
\]
Multiplying over the disjoint edge blocks gives the displayed Fourier bound.
For the concentration estimate, Fourier inversion and Lemma
\ref{lem:weighted-chain-comparison} give
\[
\begin{aligned}
        \Pp(W=z)
        &\le (2\pi)^{-r}\int_{[-\pi,\pi]^r}
        \exp\left[-c\sum_{j=1}^r A_j
        \operatorname{dist}(u_j-u_{j+1},2\pi\mathbb Z)^2\right]\dd u \\
        &\le C\prod_{j=1}^r A_j^{-1/2}.
\end{aligned}
\]
Bounded tilts only multiply the finitely many atom probabilities by bounded factors, so the edge-block lower bounds persist after changing constants.
\end{proof}

\begin{lemma}[Forward-edge projected concentration]\label{lem:projected-forward-edge-concentration}
Let $G\ne\emptyset$ be a subset of $\{1,\ldots,r\}$ and let
$0<A_1<\cdots<A_r$ satisfy $A_{j+1}/A_j\ge1+\eta$.
Let $W$ be a sum of independent random vectors supported on
$\{e_j:j\in G\}\cup\{0\}$ after projection onto the coordinates in $G$.
Suppose that for each $j\in G$ there is an index
$\kappa(j)\in (G\cap\{j+1,\ldots,r\})\cup\{0\}$ and a set of at
least $c_0A_j$ summands, disjoint as $j$ varies, for which the two
projected atoms $e_j$ and $e_{\kappa(j)}$ have probabilities at least
$c_0$; here $e_0=0$.  Then, uniformly in $z_G$,
\[
        \Pp(W_G=z_G)\le C\prod_{j\in G} A_j^{-1/2}.
\]
The same estimate holds after any exponential tilt whose coordinates are
bounded in absolute value by a fixed constant.
\end{lemma}

\begin{proof}
For a summand in the block assigned to $j$, the same two-atom argument used
in Lemma \ref{lem:edge-fourier} gives the factor
\[
        \exp\{-c\operatorname{dist}(u_j-u_{\kappa(j)},2\pi\mathbb Z)^2\},
        \qquad u_0=0 .
\]
Multiplying over the disjoint blocks gives
\[
        |\phi(u)|\le
        \exp\left[-c\sum_{j\in G}A_j
        \operatorname{dist}(u_j-u_{\kappa(j)},2\pi\mathbb Z)^2\right].
\]
Because every directed edge goes from $j$ to a larger index, or to the
anchor $0$, each path reaches $0$ after at most $r$ steps.  Along such a
path the weights only increase.  Therefore, for $u_0=0$,
\[
        \sum_{j\in G}A_j u_j^2
        \le C_{r,\eta}\sum_{j\in G}A_j(u_j-u_{\kappa(j)})^2 .
\]
The torus version follows from the same path argument.  For each $j$, choose
integer multiples of $2\pi$ successively along the directed path $P(j)$ from
$j$ to the anchor $0$.  The triangle inequality on the circle gives
\[
        \operatorname{dist}(u_j,2\pi\mathbb Z)
        \le \sum_{\ell\in P(j)}
        \operatorname{dist}(u_\ell-u_{\kappa(\ell)},2\pi\mathbb Z).
\]
Cauchy's inequality, together with monotonicity of the weights along this path,
gives the displayed comparison with distances to $2\pi\mathbb Z$.
Fourier inversion then gives
\[
        \Pp(W_G=z_G)
        \le C\int_{[-\pi,\pi]^G}
        \exp\left[-c\sum_{j\in G}A_ju_j^2\right]\dd u_G
        \le C\prod_{j\in G}A_j^{-1/2}.
\]
Bounded tilts distort all atoms in the finite support by bounded factors, so
the two-atom lower bounds persist after changing constants.
\end{proof}

\begin{lemma}[Terminal interval edge blocks]\label{lem:terminal-interval-edge-blocks}
Fix \(r\ge1\), \(0<c_*<C_*<\infty\), and \(c_0>0\).  In the terminal
interval notation above, fix \(j\in\{1,\ldots,r\}\) and suppose
\[
        c_*\le s_j=a_jt_j\le C_*,
        \qquad
        m_j(t)=\sum_{b\notin A}\Pp(T_b\in I_j)\ge c_0a_j.
\]
Then there are constants \(0<\alpha<\beta<\infty\) and \(c>0\), depending
only on \(r,c_*,C_*\), and \(c_0\), and a set of at least \(ca_j\) rates
\(b\in[\alpha a_j,\beta a_j]\setminus A\) such that
\[
        \Pp(T_b\in I_j)\ge c,
        \qquad
        \Pp(T_b\le t_j)\ge c.
\]
Consequently, after projecting any chosen subset of the interval coordinates
larger than \(j\) either to their own coordinates or to the anchor \(0\), a
further subblock of size at least \(ca_j\) gives two projected atoms
\(e_j\) and \(e_{\kappa(j)}\), where \(\kappa(j)>j\) is a retained
coordinate or \(\kappa(j)=0\), each with probability at least \(c\).
The conclusion remains valid under every interval-coordinate exponential tilt
whose coordinates are bounded by a fixed constant, after changing \(c\).
\end{lemma}

\begin{proof}
Choose \(\alpha>0\) small and \(\beta<\infty\) large.  The contribution to
\(m_j(t)\) from \(b<\alpha a_j\) is at most \(\alpha a_j+1\).  Since
\(t_j\ge c_*/a_j\), the contribution from \(b>\beta a_j\) is at most
\[
        \sum_{b>\beta a_j}e^{-bt_j}
        \le C a_j e^{-c_*\beta}.
\]
After first decreasing \(\alpha\) and then increasing \(\beta\), at least a
fixed positive multiple of \(a_j\) of the total interval mass comes from
\([\alpha a_j,\beta a_j]\).  This block contains \(O(a_j)\) rates, so at
least \(ca_j\) of them satisfy \(\Pp(T_b\in I_j)\ge c\).  For every such
rate,
\[
        bt_j\ge\alpha c_*,
        \qquad
        \Pp(T_b\le t_j)=1-e^{-bt_j}\ge1-e^{-\alpha c_*}>0.
\]
The event \(\{T_b\le t_j\}\) is the disjoint union of the later categories
\(I_{j+1},\ldots,I_{r+1}\).  For each selected rate, one of these finitely
many categories has probability bounded below.  A second pigeonhole argument
produces a common later category on a subblock of size \(ca_j\).  Projection
maps that category either to a retained atom \(e_{\kappa(j)}\) or to the anchor
\(0\).  Removing the fixed set \(A\) changes the block size by only \(O_r(1)\).
A bounded exponential tilt changes every atom probability in the finite
categorical support by a bounded multiplicative factor, which proves the final
assertion.
\end{proof}

\begin{lemma}[Projected good--bad interval bound]\label{lem:projected-good-bad-interval-bound}
Fix $r\ge1$, $\eta>0$, and a compact scaled chamber
$K\subset(0,\infty)^r$.  In the terminal interval notation above, suppose
$s=(a_1t_1,\ldots,a_rt_r)\in K$ and $t_1>\cdots>t_r$.  Suppose the deterministic targets satisfy
\[
        c_0a_j\le y_j\le C_0a_j,
        \qquad 1\le j\le r .
\]
Then, uniformly for such $t$ and $y$,
\[
        \Pp(Y(t)=y)
        \le C\prod_{j=1}^r a_j^{-1/2}
        \exp\left[-c\sum_{j=1}^r
        \min\left\{\frac{(y_j-m_j(t))^2}{a_j},a_j\right\}\right].
\]
The constants depend only on $r,\eta,K,c_0$, and $C_0$.
\end{lemma}

\begin{proof}
Write $d_j=y_j-m_j(t)$.  Since $s\in K$, finite-geometric estimates give
$m_j(t)=O_K(a_j)$ and hence $|d_j|\le M_Ka_j$.  Fix a small
$\varepsilon>0$.  Let
\[
        G=\{j: |d_j|\le\varepsilon a_j\},
        \qquad B=\{1,\ldots,r\}\setminus G .
\]
For $j\in G$ the target lower bound gives
$m_j(t)\ge(c_0-\varepsilon)a_j$; choose
$\varepsilon<c_0/2$, so Lemma \ref{lem:terminal-interval-edge-blocks}
applies to every good coordinate.

Use one combined exponential tilt.  Put
\[
        \lambda_j=
        \begin{cases}
        \delta d_j/a_j, & j\in G,\\
        \delta\,\operatorname{sgn}(d_j), & j\in B,
        \end{cases}
\]
where $\delta>0$ is a sufficiently small fixed constant.  Let
$K_t(\lambda)=\log\Ee\exp\{\lambda\cdot(Y(t)-m(t))\}$.
For every bounded tilt $\xi$, the categorical support and the finite-geometric
bound $\Var_\xi Y_j(t)\le C a_j$ give
\[
        \Var_\xi(v\cdot Y(t))\le C\sum_{j=1}^r a_jv_j^2.
\]
Therefore
\[
        K_t(\lambda)\le C\sum_{j=1}^r a_j\lambda_j^2 .
\]
On the other hand,
\[
        \lambda\cdot d
        \ge \delta\sum_{j\in G}\frac{d_j^2}{a_j}
            +\delta\varepsilon\sum_{j\in B}a_j .
\]
Choosing $\delta$ small enough yields
\[
        -\lambda\cdot d+K_t(\lambda)
        \le -c\sum_{j\in G}\frac{d_j^2}{a_j}
             -c\sum_{j\in B}a_j .
\]
Under this bounded tilt all categorical probabilities are changed by bounded
multiplicative factors.  Project the tilted array onto the good coordinates
and ignore the bad-coordinate constraints.  For each $j\in G$, Lemma \ref{lem:terminal-interval-edge-blocks} gives a candidate block $B_j$ of at least $c_0a_j$
rates with two projected atoms $e_j$ and $e_{\kappa(j)}$, where
$\kappa(j)>j$ is another good coordinate or $\kappa(j)=0$.  We now extract
disjoint subblocks rigorously.  Process the indices of $G$ in increasing order.
Since $a_{i+1}/a_i\ge1+\eta$,
\[
        \sum_{i<j}a_i\le C_{r,\eta}a_j.
\]
Choose $c'>0$ so small that
$c'(1+C_{r,\eta})\le c_0$.  After subblocks of sizes $c'a_i$ have been chosen
for all earlier good indices $i<j$, at most
$c'\sum_{i<j}a_i\le c'C_{r,\eta}a_j$ elements of $B_j$ have been used.
At least $c'a_j$ unused elements of $B_j$ therefore remain.  Continuing
greedily produces disjoint subblocks of size at least $c'a_j$ for every
$j\in G$.
Lemma \ref{lem:projected-forward-edge-concentration} gives
\[
        \Pp_\lambda(Y_G(t)=y_G)\le C\prod_{j\in G}a_j^{-1/2}.
\]
Consequently
\begin{align*}
        \Pp(Y(t)=y)
        &\le \exp\{-\lambda\cdot d+K_t(\lambda)\}
        \Pp_\lambda(Y_G(t)=y_G)  \\
        &\le C\prod_{j\in G}a_j^{-1/2}
        \exp\left[-c\sum_{j\in G}\frac{d_j^2}{a_j}
                  -c\sum_{j\in B}a_j\right].
\end{align*}
For $j\in B$, the factor $e^{-ca_j}$ absorbs the missing local factor:
$e^{-ca_j}\le C a_j^{-1/2}e^{-c'a_j}$.  Since for $j\in G$ we have
$|d_j|\le\varepsilon a_j$, this is the desired bound, after changing
constants.

\end{proof}


\begin{lemma}[Actual-mean cumulative Chernoff bounds]\label{lem:actual-mean-lower-cumulative-chernoff}
Fix \(r\ge1\).  Let \(B\subset\{1,\ldots,r\}\), and put
\[
        \mu_j=\Ee Z_j(t),\qquad j\in B.
\]
There are constants \(\delta_0,C>0\), depending only on \(r\), such that,
for every \(|\theta|\le\delta_0\),
\[
\log \Ee\exp\left\{\theta\sum_{j\in B}(Z_j(t)-\mu_j)\right\}
        \le C\theta^2\sum_{j\in B}\mu_j .
\]
After decreasing \(\delta_0\), if necessary, fix
\(0<\delta\le\delta_0\) and \(C_0<\infty\).  There are
\(c=c(\delta,r,C_0)>0\) and \(L_0=L_0(\delta,r,C_0)\) such that, whenever
\(L\ge L_0\), both of the following assertions hold.

If \(z_j\le C_0L^{-1}\mu_j\) for every \(j\in B\), then
\[
        -\sum_{j\in B}(-\delta)(z_j-\mu_j)
        +\log \Ee\exp\left\{-\delta\sum_{j\in B}(Z_j(t)-\mu_j)\right\}
        \le -c\sum_{j\in B}\mu_j .
\]
If \(\mu_j\le C_0L^{-1}z_j\) for every \(j\in B\), then
\[
        -\sum_{j\in B}\delta(z_j-\mu_j)
        +\log \Ee\exp\left\{\delta\sum_{j\in B}(Z_j(t)-\mu_j)\right\}
        \le -c\sum_{j\in B}z_j .
\]
In particular, the first exponent is at most
\(-cL\sum_{j\in B}a_j\) if \(\mu_j\ge c_0La_j\), while the second is at most
\(-c\sum_{j\in B}a_j\) if \(z_j\ge c_0a_j\).
\end{lemma}

\begin{proof}
For a single rate \(b\notin A\), define
\[
        R_b=\sum_{j\in B}\one_{\{T_b>t_j\}},
        \qquad m_b=\Ee R_b .
\]
Then \(0\le R_b\le r\), and
\[
        \sum_{b\notin A}m_b=\sum_{j\in B}\mu_j .
\]
For \(|\theta|\le1\), let
\[
        K_b(\theta)=\log \Ee e^{\theta(R_b-m_b)} .
\]
Then \(K_b(0)=K_b'(0)=0\), and
\[
        K_b''(\theta)=\Var_\theta(R_b)
        \le \Ee_\theta R_b^2
        \le r\,\Ee_\theta R_b .
\]
Since \(0\le R_b\le r\), a bounded tilt changes the expectation of the
nonnegative variable \(R_b\) by at most a factor depending only on \(r\):
\[
        \Ee_\theta R_b
        =\frac{\Ee R_be^{\theta R_b}}{\Ee e^{\theta R_b}}
        \le e^{2r}m_b,
        \qquad |\theta|\le1 .
\]
Therefore \(K_b''(\theta)\le C_rm_b\) for \(|\theta|\le1\), and Taylor's
formula gives \(K_b(\theta)\le C_r\theta^2m_b\).  Summing over
\(b\notin A\) proves the first displayed inequality for both signs of
\(\theta\).

For the lower-tail assertion, use the preceding bound and
\(z_j\le C_0L^{-1}\mu_j\):
\[
\begin{aligned}
&-\sum_{j\in B}(-\delta)(z_j-\mu_j)
 +\log \Ee\exp\left\{-\delta\sum_{j\in B}(Z_j(t)-\mu_j)\right\} \\
&\qquad\le
        -\delta\left(1-\frac{C_0}{L}\right)\sum_{j\in B}\mu_j
        +C\delta^2\sum_{j\in B}\mu_j .
\end{aligned}
\]
Decrease \(\delta_0\), if necessary, so that \(C\delta_0\le1/4\), and then
choose \(L_0\) so that \(C_0/L\le1/4\).  The last display is at most
\(-\delta\sum_{j\in B}\mu_j/2\).

For the upper-tail assertion, the same cumulant bound and
\(\mu_j\le C_0L^{-1}z_j\) give
\[
\begin{aligned}
&-\sum_{j\in B}\delta(z_j-\mu_j)
 +\log \Ee\exp\left\{\delta\sum_{j\in B}(Z_j(t)-\mu_j)\right\} \\
&\qquad\le
        -\delta\sum_{j\in B}z_j
        +(\delta+C\delta^2)\sum_{j\in B}\mu_j\\
&\qquad\le
        -\delta\left(1-\frac{2C_0}{L}\right)
        \sum_{j\in B}z_j,
\end{aligned}
\]
after decreasing \(\delta_0\) once more.  Enlarging \(L_0\) makes this at
most \(-\delta\sum_{j\in B}z_j/2\).  The two final scale bounds follow
immediately from the corresponding assumptions.
\end{proof}

\begin{lemma}[Mixed cumulative--interval localization]
\label{lem:projected-cumulative-tail-localization}
Fix $r\ge1$ and $\eta>0$.  There is $L_0=L_0(r,\eta)$ such that the
following holds for every $L\ge L_0$.  In the terminal interval notation, let
$s_j=a_jt_j$ and suppose $s\notin[L^{-1},L]^r$.  Put
\[
        B_-(s)=\{j:s_j<L^{-1}\},\qquad
        B_+(s)=\{j:s_j>L\},\qquad
        G(s)=\{1,\ldots,r\}\setminus(B_-(s)\cup B_+(s)).
\]
On every region on which these three sets are fixed,
\[
\begin{aligned}
        \Pp(Y(t)=y)
        &\le C_L\prod_{j=1}^r a_j^{-1/2}
        \exp\left[-c\sum_{j\in B_-}|H_j(t)|
                   -c(\log L)\sum_{j\in B_+}a_j\right]\\
        &\quad\times
        \exp\left[-c_L\sum_{j\in G(s)}
        \min\left\{\frac{H_j(t)^2}{a_j},a_j\right\}\right].
\end{aligned}\tag{7.3}
\]
Here $c>0$ is independent of $L$, while $C_L<\infty$ and $c_L>0$
depend only on $L,r$, and $\eta$.  Consequently,
\[
\begin{aligned}
        \Pp(Y(t)=y)
        &\le C_L\prod_{j=1}^r a_j^{-1/2}
        \exp\left[-cL\sum_{j\in B_-}a_j
                   -c(\log L)\sum_{j\in B_+}a_j\right]\\
        &\quad\times
        \exp\left[-c_L\sum_{j\in G(s)}
        a_j\min\{(s_j-s_j^*)^2,1\}\right].
\end{aligned}\tag{7.4}
\]
\end{lemma}

\begin{proof}
Throughout the proof the offender sets are fixed.  Write $B_-=B_-(s)$,
$B_+=B_+(s)$, and $G=G(s)$.  Finite-geometric estimates give, uniformly in
the admissible terminal range,
\[
\begin{array}{lll}
 j\in B_-:& \mu_j(t_j)\ge cLa_j,& H_j(t)=z_j-\mu_j(t_j)<0,\\[2mm]
 j\in B_+:& \mu_j(t_j)\le Ca_j/L,& 0<H_j(t)\le z_j\le a_j,
\end{array}
\tag{7.5}
\]
for $L$ sufficiently large.  In particular,
$|H_j(t)|\ge cLa_j$ on $B_-$ and $H_j(t)\ge ca_j$ on $B_+$.
For $j\in G$, one has $s_j\in[L^{-1},L]$, and therefore
\[
        m_j(t)\le\mu_j(t_j)\le C_La_j,
        \qquad |d_j(t)|\le M_La_j.
        \tag{7.6}
\]

Let
\[
        K_t(\theta)=\log\Ee\exp\{\theta\cdot(Y(t)-m(t))\}
\]
be the centered interval-coordinate cumulant.  If
$\gamma=(\gamma_1,\ldots,\gamma_r)$ is a cumulative-coordinate vector, put
\[
        (\mathcal T\gamma)_k=\sum_{j\ge k}\gamma_j.
\]
Since $Z_j=Y_1+\cdots+Y_j$ and $d_j=H_j-H_{j-1}$,
\[
        (\mathcal T\gamma)\cdot d=\gamma\cdot H.
        \tag{7.7}
\]
Choose $\delta>0$ so small that $3\delta$ is admissible in Lemma
\ref{lem:actual-mean-lower-cumulative-chernoff} and
$3\delta r<1/2$.  Define
\[
        \gamma_j^-=-\delta\one_{\{j\in B_-\}},
        \qquad
        \gamma_j^+=\delta(\log L)\one_{\{j\in B_+\}}.
\]

We also introduce a local interval tilt on $G$.  Ratio separation implies that
there is $c_\eta>0$ such that $y_j\ge c_\eta a_j$ for all $j$ and all
sufficiently large $n$.  Fix $0<\varepsilon<c_\eta/4$ and put
\[
        G_0=\{j\in G:|d_j|\le\varepsilon a_j\},
        \qquad G_1=G\setminus G_0.
\]
For a constant $\delta_L>0$ to be chosen below, set
\[
        \lambda_j=
        \begin{cases}
        \delta_L d_j/a_j,&j\in G_0,\\
        \delta_L\operatorname{sgn}(d_j),&j\in G_1,\\
        0,&j\notin G.
        \end{cases}
\]
Finally, let
\[
        \theta=\mathcal T\gamma^-+\mathcal T\gamma^++\lambda.
\]
Generalized H\"older inequality gives the precise combined-cumulant estimate
\[
\begin{aligned}
        K_t(\theta)
        &\le\frac13K_t(3\mathcal T\gamma^-)
             +\frac13K_t(3\mathcal T\gamma^+)
             +\frac13K_t(3\lambda).
\end{aligned}\tag{7.8}
\]
Thus, by (7.7),
\[
\begin{aligned}
-\theta\cdot d+K_t(\theta)
&\le\frac13\{-3\gamma^-\cdot H+K_t(3\mathcal T\gamma^-)\}\\
&\quad+\frac13\{-3\gamma^+\cdot H+K_t(3\mathcal T\gamma^+)\}\\
&\quad+\frac13\{-3\lambda\cdot d+K_t(3\lambda)\}.
\end{aligned}\tag{7.9}
\]
We estimate the three braces separately.

For $B_-$, (7.5) implies $z_j\le C L^{-1}\mu_j(t_j)$.  Lemma
\ref{lem:actual-mean-lower-cumulative-chernoff}, applied with parameter
$3\delta$, gives
\[
        -3\gamma^-\cdot H+K_t(3\mathcal T\gamma^-)
        \le-c\sum_{j\in B_-}\mu_j(t_j)
        \le-c\sum_{j\in B_-}|H_j(t)|.
        \tag{7.10}
\]
The same bound also retains the weaker loss
$-cL\sum_{j\in B_-}a_j$.

For $B_+$, put
\[
        R_b^+=\sum_{j\in B_+}\one_{\{T_b>t_j\}},
        \qquad 0\le R_b^+\le r.
\]
Since the summands indexed by $b\notin A$ are independent,
\[
\begin{aligned}
K_t(3\mathcal T\gamma^+)
&=\sum_{b\notin A}
 \log\Ee\exp\{3\delta(\log L)(R_b^+-\Ee R_b^+)\}\\
&\le \sum_{b\notin A}
 \log\Ee\exp\{3\delta(\log L)R_b^+\}\\
&\le L^{3\delta r}\sum_{b\notin A}\Pp(R_b^+>0)\\
&\le L^{3\delta r}\sum_{j\in B_+}\mu_j(t_j)
 \le C L^{-1/2}\sum_{j\in B_+}a_j.
\end{aligned}
\tag{7.11}
\]
On the other hand, (7.5) gives
\[
        -3\gamma^+\cdot H
        \le-c(\log L)\sum_{j\in B_+}a_j.
\]
After increasing $L_0$, (7.11) is absorbed by this linear gain, and hence
\[
        -3\gamma^+\cdot H+K_t(3\mathcal T\gamma^+)
        \le-c(\log L)\sum_{j\in B_+}a_j.
        \tag{7.12}
\]

For the local tilt, let $v$ vanish outside $G$.  Under every bounded
interval-coordinate tilt $\xi$, a categorical summand contributes at most one
coordinate, and therefore
\[
        \Var_\xi(v\cdot Y(t))
        \le\sum_{j\in G}v_j^2\Ee_\xi Y_j(t)
        \le C_L\sum_{j\in G}a_jv_j^2,
        \tag{7.13}
\]
where the last inequality follows from (7.6) and bounded distortion of atom
probabilities.  Taylor's formula for the cumulant gives
\[
        K_t(3\lambda)\le C_L\sum_{j\in G}a_j\lambda_j^2.
\]
Choose $\delta_L>0$ sufficiently small.  Then
\[
        -3\lambda\cdot d+K_t(3\lambda)
        \le-c_L\sum_{j\in G}
        \min\left\{\frac{d_j^2}{a_j},a_j\right\}.
        \tag{7.14}
\]
Indeed, on $G_0$ the linear term is a fixed multiple of $d_j^2/a_j$,
whereas on $G_1$ it is at least a fixed multiple of $a_j$.

We next recover the product local factor under the single total tilt $\theta$.
For fixed $L$, its interval coordinates are bounded by a constant depending
only on $L$ and $r$.  If $j\in G_0$, then
\[
        m_j(t)\ge y_j-\varepsilon a_j\ge(c_\eta/2)a_j.
\]
Lemma \ref{lem:terminal-interval-edge-blocks}, with
$s_j\in[L^{-1},L]$, supplies a candidate forward-edge block of order $a_j$.
Process the indices of $G_0$ increasingly and choose small subblocks greedily.
Since $a_{i+1}/a_i\ge1+\eta$ and $r$ is fixed,
\[
        \sum_{i<j}a_i\le C_{r,\eta}a_j,
\]
so disjoint subblocks of size $c_{L,r,\eta}a_j$ can be retained for every
$j\in G_0$.  Boundedness of the total tilt preserves the two-atom lower
bounds.  Lemma \ref{lem:projected-forward-edge-concentration}, applied after
projection onto $G_0$, yields
\[
        \Pp_\theta(Y_{G_0}(t)=y_{G_0})
        \le C_L\prod_{j\in G_0}a_j^{-1/2}.
        \tag{7.15}
\]
The exact change-of-measure identity and (7.9)--(7.15) now give
\[
\begin{aligned}
\Pp(Y(t)=y)
&\le C_L\prod_{j\in G_0}a_j^{-1/2}
 \exp\left[-c\sum_{j\in B_-}|H_j|
            -c(\log L)\sum_{j\in B_+}a_j\right]\\
&\quad\times
 \exp\left[-c_L\sum_{j\in G}
 \min\left\{\frac{d_j^2}{a_j},a_j\right\}\right].
\end{aligned}\tag{7.16}
\]
For $j\in G_1$, the last exponent contains $-c_La_j$; for $j\in B_-$ and
$j\in B_+$, (7.5) and the preceding tail losses contain the same fixed-order
exponential cost.  These factors absorb all missing $a_j^{-1/2}$ terms, after
changing constants.

It remains to pass from interval residuals to cumulative residuals without
losing boundary terms.  The assumptions of Lemma
\ref{lem:boundary-truncated-chain} hold with $A_j=a_j$ and $M=M_L$: (7.6)
controls $d_j$ on $G$, while (7.5) gives $|H_j|\le a_j$ on $B_+$.  Split a
fixed fraction of each offender loss in (7.16) for this comparison.  Since
$\log L$ is large, Lemma \ref{lem:boundary-truncated-chain} yields
\[
\begin{aligned}
&c_L\sum_{j\in G}\min\left\{\frac{d_j^2}{a_j},a_j\right\}
 +c\sum_{j\in B_-}|H_j|
 +c(\log L)\sum_{j\in B_+}a_j\\
&\qquad\ge
 c_L'\sum_{j\in G}\min\left\{\frac{H_j^2}{a_j},a_j\right\}
 +c'\sum_{j\in B_-}|H_j|
 +c'(\log L)\sum_{j\in B_+}a_j.
\end{aligned}
\tag{7.17}
\]
Combining (7.16) and (7.17) proves (7.3).

Finally, on $G$ the scaled coordinate $s_j$ belongs to $[L^{-1},L]$.
Uniform finite-geometric derivative estimates and strict monotonicity give
\[
        |H_j(t)|=|\mu_j(t_j^*)-\mu_j(t_j)|
        \ge c_La_j\min\{|s_j-s_j^*|,1\}.
        \tag{7.18}
\]
Together with (7.5), this converts (7.3) into (7.4).
\end{proof}


\begin{lemma}[Anisotropic tilted local upper bound]\label{lem:anisotropic-tilted-upper}
Let \(W=\sum_{\nu\in\mathcal I}W_\nu\) be as in Lemma \ref{lem:coord-fourier}.  Assume that the hypotheses of Lemma \ref{lem:coord-fourier} hold uniformly under all tilts \(\|\lambda\|_\infty\le\Lambda\).  Assume also that, for those tilted laws,
\[
        \Var_\lambda(v\cdot W)
        \le C_0\sum_{j=1}^r A_jv_j^2,
        \qquad v\in\Rr^r .
\]
Then there are constants \(C,c>0\), depending only on the constants in the assumptions, such that for every lattice point \(z\), with
\(h=z-\Ee W\),
\[
        \Pp(W=z)
        \le C\prod_{j=1}^r A_j^{-1/2}
        \exp\left[-c\sum_{j=1}^r
        \min\left\{\frac{h_j^2}{A_j}, |h_j|\right\}\right].
\]
In particular, if \(|h_j|\le A_j\) for every \(j\), then the exponent may be written as
\(-c\sum_j h_j^2/A_j\).
\end{lemma}

\begin{proof}
Let
\[
        K(\lambda)=\log\Ee e^{\lambda\cdot(W-\Ee W)}.
\]
For \(\|\lambda\|_\infty\le\Lambda\), the covariance upper bound gives
\[
        K(\lambda)
        =\int_0^1(1-s)\lambda^T
        \Var_{s\lambda}(W)\lambda\dd s
        \le C\sum_{j=1}^r A_j\lambda_j^2 .
\]
Choose a sufficiently small \(\delta>0\), depending only on the constants in the assumptions and with \(\delta\le\Lambda\), and put
\[
        \lambda_j=\delta\,\mathrm{sgn}(h_j)
        \min\{|h_j|/A_j,1\}.
\]
By tilting and then applying Lemma \ref{lem:coord-fourier} under the tilted law,
\[
\begin{aligned}
        \Pp(W=z)
        &=\exp\{-\lambda\cdot h+K(\lambda)\}
          \Pp_\lambda(W=z) \\
        &\le C\prod_{j=1}^r A_j^{-1/2}
        \exp\{-\lambda\cdot h+K(\lambda)\}.
\end{aligned}
\]
For the chosen \(\lambda\),
\[
        \lambda\cdot h
        =\delta\sum_{j=1}^r |h_j|\min\{|h_j|/A_j,1\},
\]
while
\[
        K(\lambda)
        \le C\delta^2\sum_{j=1}^r
        A_j\min\{h_j^2/A_j^2,1\}.
\]
Taking \(\delta\) small enough makes the second display at most one half of the first, after changing constants.  This proves the claim.
\end{proof}

\begin{lemma}[Anisotropic terminal local CLT]\label{lem:terminal-llt}
Fix \(r\ge1\), \(\eta>0\), \(\kappa>0\), and \(0<c_0<C_0<\infty\).  Let
\[
        1\ll a_1<\cdots<a_r\ll n,
        \qquad a_{j+1}/a_j\ge1+\eta,
\]
where the ratio-separation condition is empty when \(r=1\).  Let
\(t_1>\cdots>t_r\) satisfy
\[
        \frac{c_0}{a_j}\le t_j\le \frac{C_0}{a_j},
        \qquad
        \frac{t_j}{t_{j+1}}\ge1+\kappa
        \quad (1\le j<r).
\]
Let \(A=\{a_1,\ldots,a_r\}\), and put
\[
        Z_j(t)=\sum_{b\notin A}\one_{\{T_b>t_j\}},
        \qquad Z(t)=(Z_1(t),\ldots,Z_r(t)).
\]
If \(\Sigma(t)=\Var Z(t)\), then there are constants \(c,C>0\), depending only on \(r,\eta,\kappa,c_0,C_0\), such that
\[
        c\,\diag(a_1,\ldots,a_r)
        \le \Sigma(t)\le
        C\,\diag(a_1,\ldots,a_r)
\]
in quadratic-form order.  Moreover,
\[
\begin{aligned}
        \Pp(Z(t)=z)
        &=(2\pi)^{-r/2}(\det\Sigma(t))^{-1/2} \\
        &\quad\times
        \exp\left[-\frac12(z-\Ee Z(t))^T\Sigma(t)^{-1}(z-\Ee Z(t))\right]
        +o\left((a_1\cdots a_r)^{-1/2}\right)
\end{aligned}
\]
uniformly whenever the Mahalanobis norm of \(z-\Ee Z(t)\) is bounded.  In addition, for every fixed \(\gamma\in(0,1/6)\), the coordinatewise moderate-deviation upper bound
\[
        \Pp(Z(t)=z)
        \le C(\det\Sigma(t))^{-1/2}
        \exp\left[-c\sum_{j=1}^r
        \frac{(z_j-\Ee Z_j(t))^2}{a_j}\right]
\]
holds uniformly whenever
\[
        |z_j-\Ee Z_j(t)|\le a_j^{1/2+\gamma},
        \qquad 1\le j\le r .
\]
\end{lemma}

\begin{proof}
For a single rate \(b\notin A\), write
\[
        V_b(t)=(\one_{\{T_b>t_1\}},\ldots,\one_{\{T_b>t_r\}}).
\]
The vector \(V_b(t)\) takes the cumulative values
\[
        (1,\ldots,1),\quad (0,1,\ldots,1),\quad\ldots,
        \quad (0,\ldots,0,1),\quad (0,\ldots,0),
\]
according as \(T_b\) falls in
\[
        (t_1,\infty),\quad (t_2,t_1],\quad\ldots,
        \quad(t_r,t_{r-1}],\quad (0,t_r].
\]
Adjacent support points differ by the coordinate vectors, so the lattice span is \(\mathbb Z^r\).

Choose constants \(0<d_1<d_2<\infty\), depending only on
\(r,\eta,\kappa,c_0,C_0\), with \(d_2/d_1<1+\eta\).  Then the rate intervals
\([d_1a_j,d_2a_j]\) are disjoint for different \(j\)'s.  For every integer rate \(b\in[d_1a_j,d_2a_j]\) and all large \(n\), the two adjacent time categories separated by the threshold \(t_j\) have probabilities bounded below.  For instance,
\[
        e^{-bt_j}\ge e^{-d_2C_0},
        \qquad
        b(t_{j-1}-t_j)\ge d_1c_0\frac{\kappa}{1+\kappa}
\]
for the upper adjacent interval, with the evident endpoint modifications when \(j=1\) or \(j=r\).  The fixed omitted set \(A\) removes only \(r\) rates, so each block still contains order \(a_j\) available rates.

These disjoint adjacent blocks imply the covariance lower bound: for every \(u\in\Rr^r\), the \(j\)-th block contributes at least \(c a_j u_j^2\) to \(\Var(u\cdot Z(t))\).  Thus
\[
        \Var(u\cdot Z(t))\ge c\sum_{j=1}^r a_j u_j^2 .
\]
The upper bound follows from \(\Var Z_j(t)=O(a_j)\), which is a direct finite-geometric estimate, and Cauchy's inequality:
\[
        \Var(u\cdot Z(t))
        \le r\sum_{j=1}^r u_j^2\Var Z_j(t)
        \le C\sum_{j=1}^r a_j u_j^2 .
\]
This proves the covariance comparison.

The same blocks give the coordinatewise Fourier bound of Lemma \ref{lem:coord-fourier}, with \(A_j=a_j\).  This is the point at which anisotropy is used: rather than bounding the whole off-origin region by a scalar factor \(e^{-ca_1}\), we retain the product-scale estimate
\[
        |\Ee e^{iu\cdot Z(t)}|
        \le \exp\left[-c\sum_{j=1}^r
        a_j\operatorname{dist}(u_j,2\pi\mathbb Z)^2\right].
\]
This bound is stable under bounded exponential tilts, because the support is finite and the two distinguished adjacent atom probabilities remain bounded below.

We next prove the local CLT.  Fourier inversion gives
\[
        \Pp(Z(t)=z)
        =(2\pi)^{-r}\int_{[-\pi,\pi]^r}
        e^{-iu\cdot(z-\Ee Z(t))}\phi_t(u)\dd u,
\]
where \(\phi_t(u)=\Ee e^{iu\cdot(Z(t)-\Ee Z(t))}\).  Put
\(D_a=\diag(a_1,\ldots,a_r)\).  On the anisotropic small box
\(|u_j|\le R a_j^{-1/2}\), the finite-tail estimates and boundedness of the increments give the uniform expansion
\[
        \log\phi_t(u)
        =-\frac12u^T\Sigma(t)u
        +O\left((\max_j |u_j|)\sum_{j=1}^r a_j u_j^2\right).
\]
Indeed, the third cumulant in the direction \(u\) is bounded by a constant times \(\max_j|u_j|\) times the corresponding variance contribution, and the variance upper bound just proved controls the latter by \(C\sum_j a_j u_j^2\).  After the change of variables \(u=D_a^{-1/2}v\), the error term is
\(O_R(a_1^{-1/2})=o(1)\).  Since the Mahalanobis norm of
\(z-\Ee Z(t)\) is bounded, this gives the Gaussian integral on every fixed \(v\)-ball.  The contribution from the complement of such a ball is controlled uniformly by the coordinatewise Fourier bound above:
\[
        \int_{\max_j |u_j|>R a_j^{-1/2}}
        \exp\left[-c\sum_j a_j\operatorname{dist}(u_j,2\pi\mathbb Z)^2\right]\dd u
        \le C\left(\prod_j a_j^{-1/2}\right)e^{-cR^2}.
\]
Letting \(R\to\infty\) yields the stated local CLT and the error
\(o((a_1\cdots a_r)^{-1/2})\), uniformly over the allowed parameters.

For the moderate-deviation upper bound, apply Lemma \ref{lem:anisotropic-tilted-upper} to the terminal array.  The covariance upper bound required there has already been proved and remains valid under the small tilts used below.  If
\(|z_j-\Ee Z_j(t)|\le a_j^{1/2+\gamma}\), then for large \(n\) this displacement is at most \(a_j\), and Lemma \ref{lem:anisotropic-tilted-upper} gives
\[
        \Pp(Z(t)=z)
        \le C\prod_{j=1}^r a_j^{-1/2}
        \exp\left[-c\sum_{j=1}^r
        \frac{(z_j-\Ee Z_j(t))^2}{a_j}\right].
\]
The covariance comparison implies
\(\prod_j a_j^{-1/2}\asymp(\det\Sigma(t))^{-1/2}\), proving the displayed form.
\end{proof}


\begin{lemma}[Scaled interval-count localization]
\label{lem:terminal-interval-tilt}
Fix $r\ge1$, $\eta>0$, and a compact set $K\subset(0,\infty)^r$.  In the
terminal interval notation, define
\[
        \Phi(t)=\prod_{j=1}^r a_j e^{-a_jt_j}.
\]
Uniformly for $s=(s_1,\ldots,s_r)\in K$ with $t_1>\cdots>t_r$,
\[
        \Pp(Y(t)=y)
        \le C\left(\prod_{j=1}^r a_j^{-1/2}\right)
        \exp\left[-c\sum_{j=1}^r
        a_j\min\{(s_j-s_j^*)^2,1\}\right],
        \tag{7.19}
\]
where $C,c>0$ depend only on $r,\eta$, and $K$.

For a precise uniform tail statement, let $\mathfrak A_N(\eta)$ be the set of
integer tuples $(n,a_1,\ldots,a_r)$ satisfying
\[
        n\ge N,\qquad
        N\le a_1<\cdots<a_r\le n/N,
        \qquad a_{j+1}/a_j\ge1+\eta.
\]
Then
\[
\begin{aligned}
\lim_{L\to\infty}\limsup_{N\to\infty}
\sup_{(n,a)\in\mathfrak A_N(\eta)}
&\frac{1}{\prod_{j=1}^r a_j^{-1}}
\int_{\substack{t_1>\cdots>t_r\\
        (a_1t_1,\ldots,a_rt_r)\notin[L^{-1},L]^r}}
        \Phi(t)\Pp(Z(t)=z)\dd t
=0.
\end{aligned}\tag{7.20}
\]
In particular, for every $\varepsilon>0$ there are $L<\infty$ and $N_0$ such
that the integral in (7.20) is at most
$\varepsilon\prod_j a_j^{-1}$ for every
$(n,a)\in\mathfrak A_{N_0}(\eta)$.
\end{lemma}

\begin{proof}
Ratio separation gives
\[
        c_\eta a_j\le y_j\le C_\eta a_j,
        \qquad 1\le j\le r,
\]
for all sufficiently large $a_1$.  On the fixed compact scaled chamber $K$,
finite-geometric estimates give $m_j(t)=O_K(a_j)$ and hence
$|d_j(t)|\le M_Ka_j$.  Lemma
\ref{lem:projected-good-bad-interval-bound} therefore yields
\[
        \Pp(Y(t)=y)
        \le C\prod_{j=1}^r a_j^{-1/2}
        \exp\left[-c\sum_{j=1}^r
        \min\left\{\frac{d_j(t)^2}{a_j},a_j\right\}\right].
\]
Since $d_j=H_j-H_{j-1}$ and $H_0=0$, Lemma
\ref{lem:truncated-triangular-comparison} gives
\[
        \sum_{j=1}^r
        \min\left\{\frac{d_j(t)^2}{a_j},a_j\right\}
        \ge c_K\sum_{j=1}^r
        \min\left\{\frac{H_j(t)^2}{a_j},a_j\right\}.
\]
On $K$, the derivative of
$s\mapsto\mu_j(s/a_j)$ has magnitude comparable to $a_j$, uniformly in $j$.
The mean-value theorem near the saddle, followed by compactness and strict
monotonicity away from it, gives
\[
        |H_j(t)|
        \ge c_Ka_j\min\{|s_j-s_j^*|,1\}.
\]
This proves (7.19).

We turn to (7.20).  Fix $L\ge L_0(r,\eta)$ and partition the complement of
$[L^{-1},L]^r$ according to the finitely many possible pairs of offender sets
$(B_-,B_+)$.  On each such region, Lemma
\ref{lem:projected-cumulative-tail-localization} gives (7.4).  Since
\[
        \Phi(t)\dd t_1\cdots\dd t_r
        =\prod_{j=1}^r e^{-s_j}\dd s_j,
\]
and discarding the chamber constraint only enlarges the integration domain,
the resulting bound factors coordinatewise.  Every nonoffending coordinate
contributes at most
\[
        C_L\int_{L^{-1}}^L
        e^{-c_La_j\min\{(s-s_j^*)^2,1\}}\dd s
        \le C_La_j^{-1/2}.
        \tag{7.21}
\]
For a lower offender, the integration interval has length at most $L^{-1}$,
and
\[
        a_j^{1/2}e^{-cLa_j}
        \le \alpha_{L,N}^-
        :=\sup_{a\ge N}a^{1/2}e^{-cLa}
        \longrightarrow0
        \qquad(N\to\infty).
        \tag{7.22}
\]
For an upper offender,
\[
        \int_L^\infty e^{-s}\dd s=e^{-L},
\]
and
\[
        a_j^{1/2}L^{-ca_j}
        \le \alpha_{L,N}^+
        :=\sup_{a\ge N}a^{1/2}L^{-ca}
        \longrightarrow0
        \qquad(N\to\infty).
        \tag{7.23}
\]
The original factor $\prod_j a_j^{-1/2}$, together with (7.21)--(7.23),
therefore supplies a factor $a_j^{-1}$ for every coordinate and an additional
quantity tending to zero for every nonempty offender set.  The number of
possible offender sets is finite.  Hence, for each fixed $L\ge L_0$,
\[
\limsup_{N\to\infty}
\sup_{(n,a)\in\mathfrak A_N(\eta)}
\frac{
\displaystyle
\int_{\substack{t_1>\cdots>t_r\\s\notin[L^{-1},L]^r}}
\Phi(t)\Pp(Z(t)=z)\dd t}
{\prod_j a_j^{-1}}=0.
\]
This implies (7.20) and the final assertion.
\end{proof}


\begin{lemma}[Product-scale terminal localization]\label{lem:terminal-product-localization}
Fix \(r\ge1\) and \(\eta>0\).  In the setting of Lemma \ref{lem:suk-joint-saddle}, put
\[
        z=(a_1-1,a_2-2,\ldots,a_r-r),
        \qquad
        \Phi(t)=\prod_{j=1}^r a_j e^{-a_jt_j},
\]
and let \(Z(t)\) have the same definition as in Lemma \ref{lem:terminal-llt}.  There is \(\delta_0>0\), depending only on \(r\) and \(\eta\), such that for every fixed \(0<\delta<\delta_0\) there are constants \(C,c>0\), depending only on \(r,\eta,\delta\), for which the following holds.  If
\[
        \mathcal N_\delta
        =\left\{t_1>\cdots>t_r:
        |t_j/t_j^*-1|\le\delta\text{ for all }j\right\},
\]
and
\[
        u_j=a_j^{3/2}(t_j-t_j^*),
        \qquad 1\le j\le r,
\]
then, uniformly for \(t\in\mathcal N_\delta\),
\[
        \Phi(t)\Pp(Z(t)=z)\dd t_1\cdots\dd t_r
        \le C\left(\prod_{j=1}^r a_j^{-1}\right)
        \exp\left[-c\sum_{j=1}^r\min(u_j^2,a_j)\right]
        \dd u_1\cdots\dd u_r .
\]
Moreover,
\[
\begin{aligned}
&\int_{\{t_1>\cdots>t_r\}\setminus\mathcal N_\delta}
        \Phi(t)\Pp(Z(t)=z)\dd t_1\cdots\dd t_r  \\
&\qquad =o\left(\prod_{j=1}^r a_j^{-1}\right),
\end{aligned}
\]
uniformly over the log-separated terminal range.
\end{lemma}

\begin{proof}
Choose \(\delta_0\) small enough that, whenever
\(|t_j/t_j^*-1|\le\delta_0\) for all \(j\), the ratios
\(t_j/t_{j+1}\) are bounded below by \(1+\kappa\), for some
\(\kappa>0\) depending only on \(\eta\).  On this neighborhood, Lemma \ref{lem:suk-joint-saddle} gives
\[
        t_j\asymp a_j^{-1},\qquad
        -\mu_j'(t_j)\asymp a_j^2,
        \qquad
        \mu_j''(t_j)=O(a_j^3),
\]
uniformly.

For \(t\in\mathcal N_\delta\), put \(s_j=a_jt_j\).  Since
\[
        s_j-s_j^*=a_j(t_j-t_j^*)=u_j/\sqrt{a_j},
\]
Lemma \ref{lem:terminal-interval-tilt}, applied on the compact scaled neighborhood determined by \(\delta\), gives
\[
        \Pp(Z(t)=z)
        \le C\left(\prod_{j=1}^r a_j^{-1/2}\right)
        \exp\left[-c\sum_{j=1}^r\min(u_j^2,a_j)\right].
\]
Also
\[
        \Phi(t)\dd t_1\cdots\dd t_r
        =\left(\prod_{j=1}^r a_j e^{-s_j}\right)
        \left(\prod_{j=1}^r a_j^{-3/2}\right)
        \dd u_1\cdots\dd u_r
        \le C\prod_{j=1}^r a_j^{-1/2}\dd u_1\cdots\dd u_r,
\]
because \(s_j\) remains in a compact subset of \((0,\infty)\).  Multiplying the two displays proves the product-scale bound.

It remains to remove the complement of \(\mathcal N_\delta\).  In the scaled variables, Lemma \ref{lem:terminal-interval-tilt} first removes the complement of a fixed compact set at cost
\(o(\prod_j a_j^{-1})\).  On the remaining compact set, the same lemma gives
\[
        \Pp(Z(t)=z)
        \le C\left(\prod_{j=1}^r a_j^{-1/2}\right)
        \exp\left[-c\sum_{j=1}^r
        a_j\min\{(s_j-s_j^*)^2,1\}\right].
\]
If \(t\notin\mathcal N_\delta\), then, on this compact set and for all large \(n\), at least one scaled coordinate satisfies
\(|s_j-s_j^*|\ge c_\delta\).  Since
\[
        \Phi(t)\dd t_1\cdots\dd t_r
        =\prod_{j=1}^r e^{-s_j}\dd s_j,
\]
the compact contribution outside \(\mathcal N_\delta\) is bounded by
\[
        C\left(\prod_{j=1}^r a_j^{-1/2}\right)
        \int_{\max_j|s_j-s_j^*|\ge c_\delta}
        \exp\left[-c\sum_{j=1}^r
        a_j\min\{(s_j-s_j^*)^2,1\}\right]\dd s .
\]
Changing variables \(v_j=\sqrt{a_j}(s_j-s_j^*)\) gives
\[
        O\left(e^{-c_\delta a_1}\prod_{j=1}^r a_j^{-1}\right)
        =o\left(\prod_{j=1}^r a_j^{-1}\right).
\]
Together with the compact-tail estimate, this proves the second assertion.
\end{proof}

\begin{lemma}[Covariance stability on the terminal saddle scale]
\label{lem:terminal-covariance-stability}
Fix \(r\ge1\), \(\eta>0\), and \(R<\infty\).  In the setting of Lemma
\ref{lem:suk-joint-saddle}, let
\[
        \mathsf D_a=\diag(a_1,\ldots,a_r),
        \qquad
        \mathsf H_a=\diag(a_1^{-3/2},\ldots,a_r^{-3/2}),
\]
and let \(\Sigma(t)=\Var Z(t)\).  Then, uniformly over the log-separated
terminal range,
\[
\sup_{\|u\|\le R}
\left\|
\mathsf D_a^{-1/2}
\bigl(\Sigma(t^*+\mathsf H_au)-\Sigma(t^*)\bigr)
\mathsf D_a^{-1/2}
\right\|_{\mathrm{op}}
\longrightarrow0.
\tag{7.24}
\]
Consequently, uniformly on the same fixed \(u\)-ball,
\[
        \frac{\det\Sigma(t^*+\mathsf H_au)}{\det\Sigma(t^*)}
        \longrightarrow1
\tag{7.25}
\]
and
\[
\left\|
\mathsf D_a^{1/2}
\bigl(\Sigma(t^*+\mathsf H_au)^{-1}-\Sigma(t^*)^{-1}\bigr)
\mathsf D_a^{1/2}
\right\|_{\mathrm{op}}
\longrightarrow0.
\tag{7.26}
\]
\end{lemma}

\begin{proof}
For \(j\le k\), independence of the nonspecified clocks gives
\[
        \Sigma_{jk}(t)
        =\sum_{b\notin A}e^{-bt_j}(1-e^{-bt_k}).
        \tag{7.27}
\]
The same formula with \(j=k\) is the Bernoulli variance.  On a fixed saddle
ball \(t=t^*+\mathsf H_au\), Lemma \ref{lem:suk-joint-saddle} gives
\(t_j\asymp a_j^{-1}\), uniformly.  Differentiating (7.27) and using the
finite-geometric estimates yields
\[
        |\partial_{t_j}\Sigma_{jk}(t)|
        \le\sum_{b\ge1}b e^{-bt_j}\le C_Ra_j^2,
\]
while
\[
        |\partial_{t_k}\Sigma_{jk}(t)|
        \le\sum_{b\ge1}b e^{-b(t_j+t_k)}\le C_Ra_j^2.
\]
For \(j=k\), the derivative is bounded by the same geometric sum.  Hence
\[
\begin{aligned}
|\Sigma_{jk}(t^*+\mathsf H_au)-\Sigma_{jk}(t^*)|
&\le C_Ra_j^2(a_j^{-3/2}+a_k^{-3/2}).
\end{aligned}
\]
After division by \(\sqrt{a_ja_k}\), the right-hand side is at most
\[
        C_R\left(a_k^{-1/2}+\frac{a_j^{3/2}}{a_k^2}\right)
        \le C_Ra_1^{-1/2}.
\]
Since \(r\) is fixed, entrywise convergence implies (7.24).

Lemma \ref{lem:terminal-llt} gives uniform positive lower and upper bounds for
the eigenvalues of the normalized matrices
\[
        \mathsf D_a^{-1/2}\Sigma(t)\mathsf D_a^{-1/2}
\]
throughout the saddle neighborhood.  The determinant and inverse are
continuous, uniformly Lipschitz functions on this compact set of positive
definite matrices.  Applying these facts to (7.24) proves (7.25) and (7.26).
\end{proof}

\begin{lemma}[Terminal integrated saddle]\label{lem:terminal-integrated-saddle}
Fix \(r\ge1\) and \(\eta>0\).  In the setting of Lemma \ref{lem:suk-joint-saddle}, put
\[
        z=(a_1-1,a_2-2,\ldots,a_r-r),
        \qquad
        \Phi(t)=\prod_{j=1}^r a_j e^{-a_jt_j}.
\]
Let \(A=\{a_1,\ldots,a_r\}\), and put
\[
        Z_j(t)=\sum_{b\notin A}\one_{\{T_b>t_j\}},
        \qquad Z(t)=(Z_1(t),\ldots,Z_r(t)).
\]
Then
\[
\begin{aligned}
&\int_{t_1>\cdots>t_r}
        \Phi(t)\Pp(Z(t)=z)\dd t_1\cdots\dd t_r  \\
&\qquad =
        (1+o(1))\Phi(t^*)
        \prod_{j=1}^r\bigl(-\mu_j'(t_j^*)\bigr)^{-1},
\end{aligned}
\]
uniformly over the log-separated terminal range.
\end{lemma}

\begin{proof}
Let \(\mathcal N_\delta\) be the product saddle neighborhood from Lemma \ref{lem:terminal-product-localization}, with \(\delta>0\) fixed and sufficiently small.  That lemma gives
\[
        \int_{\{t_1>\cdots>t_r\}\setminus\mathcal N_\delta}
        \Phi(t)\Pp(Z(t)=z)\dd t_1\cdots\dd t_r
        =o\left(\prod_{j=1}^r a_j^{-1}\right).
\]
Since Lemma \ref{lem:suk-joint-saddle} gives
\[
        \Phi(t^*)\prod_{j=1}^r(-\mu_j'(t_j^*))^{-1}
        \asymp \prod_{j=1}^r a_j^{-1},
\]
the contribution outside \(\mathcal N_\delta\) is negligible on the required relative scale.

Inside \(\mathcal N_\delta\), put
\[
        H=\diag(a_1^{-3/2},\ldots,a_r^{-3/2}),
        \qquad
        t=t^*+Hu .
\]
This scale is forced by the balance
\(-\mu_j'(t_j^*)\asymp a_j^2\) and \(\Var Z_j(t^*)\asymp a_j\): changing \(t_j\) by \(a_j^{-3/2}\) moves the mean by order \(a_j^{1/2}\), the natural standard-deviation scale.  Let
\[
        D=\diag\bigl(-\mu_1'(t_1^*),\ldots,-\mu_r'(t_r^*)\bigr).
\]
For bounded \(u\), Lemma \ref{lem:suk-joint-saddle} and Taylor's theorem give the normalized mean expansion
\[
        \left\|
        \mathsf D_a^{-1/2}
        \bigl(\Ee Z(t)-z+DHu\bigr)
        \right\|\longrightarrow0,
        \qquad \mathsf D_a=\diag(a_1,\ldots,a_r),
        \tag{7.28}
\]
uniformly on every fixed \(u\)-ball.  Indeed, the \(j\)-th Taylor remainder is
\[
        O(a_j^3|t_j-t_j^*|^2)=O(u_j^2)=o(a_j^{1/2}),
\]
whereas
\[
        \mu_j'(t_j^*)(t_j-t_j^*)=-u_j a_j^{1/2}(1+o(1)).
\]
The neighborhood \(\mathcal N_\delta\) supplies constants \(\kappa,c_0,C_0\) for the hypotheses of Lemma \ref{lem:terminal-llt}; hence the anisotropic local CLT applies uniformly for bounded \(u\).  Lemma
\ref{lem:terminal-covariance-stability}, together with (7.28), gives
\[
\begin{aligned}
        (\det\Sigma(t))^{-1/2}
        &=(\det\Sigma(t^*))^{-1/2}(1+o(1)),\\
        (z-\Ee Z(t))^T\Sigma(t)^{-1}(z-\Ee Z(t))
        &=u^TH^TD^T\Sigma(t^*)^{-1}DHu+o(1),
\end{aligned}
        \tag{7.29}
\]
uniformly on every fixed ball.  Also,
\[
        \Phi(t^*+Hu)=\Phi(t^*)(1+o(1))
\]
uniformly there, because \(a_j(t_j-t_j^*)=u_j/\sqrt{a_j}=o(1)\).
Consequently, on \(\|u\|\le R\), the integrand has the uniform asymptotic
Gaussian form
\[
        \Phi(t^*)(2\pi)^{-r/2}(\det\Sigma(t^*))^{-1/2}
        \exp\left[-\frac12 u^T H^TD^T\Sigma(t^*)^{-1}DHu\right]
        |\det H|\,\dd u .
\]
The matrix \(H^TD^T\Sigma(t^*)^{-1}DH\) is uniformly positive definite and uniformly bounded in quadratic-form order, because \(D_{jj}\asymp a_j^2\), \(H_{jj}=a_j^{-3/2}\), and \(\Sigma(t^*)\asymp\diag(a_1,\ldots,a_r)\).  Hence, as \(n\to\infty\) with \(R\) fixed, the contribution from \(\|u\|\le R\) is
\[
\begin{aligned}
        &\Phi(t^*)(2\pi)^{-r/2}(\det\Sigma(t^*))^{-1/2}|\det H| \\
        &\quad\times
        \int_{\|u\|\le R}
        \exp\left[-\frac12 u^TH^TD^T\Sigma(t^*)^{-1}DHu\right]
        \dd u
        +o_R\left(\prod_{j=1}^r a_j^{-1}\right).
\end{aligned}
\]
Letting \(R\to\infty\), uniformly over the allowed rate tuples, the Gaussian integral gives the determinant cancellation
\[
\begin{aligned}
&(2\pi)^{-r/2}(\det\Sigma(t^*))^{-1/2}|\det H|
 \int_{\mathbb R^r}
        \exp\left[-\frac12 u^TH^TD^T\Sigma(t^*)^{-1}DHu\right]\dd u  \\
&\qquad = |\det D|^{-1}.
\end{aligned}
\]
Indeed,
\[
        \det(H^TD^T\Sigma(t^*)^{-1}DH)
        =|\det H|^2|\det D|^2(\det\Sigma(t^*))^{-1}.
\]
Thus the fixed-ball saddle contribution tends to
\[
        \Phi(t^*)\prod_{j=1}^r(-\mu_j'(t_j^*))^{-1}.
\]

It remains only to justify the passage from fixed \(R\) to the whole neighborhood.  This is precisely the product-scale bound from Lemma \ref{lem:terminal-product-localization}: on \(\mathcal N_\delta\),
\[
        \Phi(t)\Pp(Z(t)=z)\dd t
        \le C\left(\prod_{j=1}^r a_j^{-1}\right)
        \exp\left[-c\sum_{j=1}^r\min(u_j^2,a_j)\right]\dd u .
\]
The functions on the right have uniformly integrable tails, since
\[
        \sup_{a\ge1}\int_{\mathbb R}e^{-c\min(u^2,a)}
        \one_{\{|u|>R\}}\one_{\{|u|\le C\sqrt a\}}\dd u\to0
        \qquad (R\to\infty).
\]
Consequently the contribution from \(\mathcal N_\delta\cap\{\|u\|>R\}\) is
\[
        o_R\left(\prod_{j=1}^r a_j^{-1}\right)
\]
uniformly as \(R\to\infty\).  Combining this tail estimate with the negligible contribution outside \(\mathcal N_\delta\) proves the lemma.
\end{proof}

\begin{proposition}[Separated terminal joint estimate]\label{prop:suk-joint}
Fix \(r\ge1\) and \(\eta>0\).  Uniformly over integers
\[
        1\ll a_1<\cdots<a_r\ll n
\]
satisfying
\[
        a_{j+1}/a_j\ge 1+\eta \quad (1\le j<r),
\]
one has
\[
        \Pp(I_{n,a_1}=\cdots=I_{n,a_r}=1)
        \sim \frac{e^{-r}}{a_1\cdots a_r}.
\]
\end{proposition}

\begin{proof}
Let \(A=\{a_1,\ldots,a_r\}\).  On the event that all \(I_{n,a_j}=1\) occur, the time ordering is forced.  Rate \(a\) is deleted at the transition at which the number of alive particles drops from \(a\) to \(a-1\).  Since \(a_1<\cdots<a_r\), the deletion at level \(a_r\) occurs before the deletion at level \(a_{r-1}\), and so on.  Thus necessarily
\[
        T_{a_1}>T_{a_2}>\cdots>T_{a_r}.
\]
Condition on \(T_{a_j}=t_j\) in this chamber.  For \(b\notin A\), define
\[
        Z_j(t)=\sum_{b\notin A}\one_{\{T_b>t_j\}},
        \qquad 1\le j\le r.
\]
At time \(t_j\), precisely \(j-1\) of the specified clocks have larger times, namely \(a_1,\ldots,a_{j-1}\).  Therefore the joint event is equivalent to
\[
        Z_j(t)=a_j-j,
        \qquad 1\le j\le r.
\]
Conversely, in the chamber \(t_1>\cdots>t_r\), these equations imply that each specified rate \(a_j\) has exactly \(a_j-1\) clocks with larger time, and hence all the events \(I_{n,a_j}=1\) occur.  Consequently the joint probability is exactly
\[
\begin{aligned}
\int_{t_1>\cdots>t_r}
        \left(\prod_{j=1}^r a_j e^{-a_j t_j}\right)
        \Pp(Z_j(t)=a_j-j,
        1\le j\le r)
        \dd t_1\cdots\dd t_r .
\end{aligned}
\]
By Lemma \ref{lem:terminal-integrated-saddle}, this integral equals
\[
        (1+o(1))
        \prod_{j=1}^r
        \frac{a_j e^{-a_jt_j^*}}{-\mu_j'(t_j^*)}.
\]
Lemma \ref{lem:suk-joint-saddle} gives \(a_jt_j^*\to1\) and \(-\mu_j'(t_j^*)\sim a_j^2\).  Hence
\[
        \prod_{j=1}^r
        \frac{a_j e^{-a_jt_j^*}}{-\mu_j'(t_j^*)}
        \sim
        \prod_{j=1}^r\frac{a_je^{-1}}{a_j^2}
        =\frac{e^{-r}}{a_1\cdots a_r}.
\]
\end{proof}

\subsection{Terminal point process and full factorial moments}

\begin{theorem}[Local terminal Poisson point process]\label{thm:suk-ppp}
Let \(b_n\to\infty\) and \(b_n=o(n)\), and define
\[
        \Xi_n=\sum_{a=1}^n I_{n,a}\delta_{\log(a/b_n)}.
\]
For every bounded interval \(J\subset\mathbb R\),
\[
        \Xi_n|_J\Rightarrow \PPP(e^{-1}\dd x)|_J
\]
in the space of locally finite counting measures on \(J\), equipped with the
vague topology.  Equivalently, counts in finitely many disjoint bounded
intervals converge jointly to independent Poisson random variables with means
\(e^{-1}\) times their lengths.
\end{theorem}

\begin{proof}
It suffices to prove convergence of mixed factorial moments.  If \(\log(a/b_n)\in J\), then \(a\to\infty\), \(a=o(n)\), and all such \(a\)'s are within a bounded multiplicative factor.  Also
\[
        \sum_{\log(a/b_n)\in J}\frac1a\to |J|.
\]
For fixed \(r\), expand \(\Ee(\Xi_n(J))_r\) over ordered distinct \(r\)-tuples.  First restrict to tuples satisfying \(|\log a_i-\log a_j|>\eta\) for all \(i\ne j\).  After sorting such a tuple increasingly, the separation condition implies
\[
        \frac{a_{j+1}}{a_j}\ge e^\eta=1+(e^\eta-1),
\]
so Proposition \ref{prop:suk-joint} applies uniformly, with separation parameter \(e^\eta-1\), and gives
\[
        e^{-r}\sum_{\eta\text{-sep}}^*\frac1{a_1\cdots a_r}+o(1).
\]
As \(n\to\infty\), this tends to \(e^{-r}\) times the Lebesgue measure of the \(\eta\)-separated part of \(J^r\).  Letting \(\eta\downarrow0\) gives \((e^{-1}|J|)^r\).

The excluded tuples are controlled by Lemma \ref{lem:death-bound}.  Indeed,
\[
\begin{aligned}
&\sum_{\substack{\log(a_i/b_n)\in J\,\,1\le i\le r\\
        |\log a_i-\log a_j|\le\eta\text{ for some }i\ne j}}
        \prod_{i=1}^r\frac1{a_i} \\
&\qquad\longrightarrow
        \operatorname{Leb}\{(x_1,\ldots,x_r)\in J^r:
        |x_i-x_j|\le\eta\text{ for some }i\ne j\}
        =O_{r,J}(\eta).
\end{aligned}
\]
Thus the excluded contribution vanishes as \(\eta\downarrow0\).  For finitely many disjoint intervals \(J_1,\ldots,J_m\), the mixed factorial moment expansion is the same expansion with each selected rate colored by the interval from which it comes.  The separated part factors into the corresponding products of lengths, and the same non-separated estimate removes the diagonals.  These limits are the mixed factorial moments of independent Poisson variables
with means \(e^{-1}|J_1|,\ldots,e^{-1}|J_m|\), so the corresponding count
vectors converge in distribution.

It remains to pass from interval counts to convergence of point measures.  The
death-process bound gives, for every fixed bounded interval \(J\),
\[
        \sup_n\Ee\Xi_n(J)
        \le C_J<\infty.
        \tag{7.30}
\]
Indeed, the relevant terminal coordinates lie in a bounded multiplicative
window and
\(\Pp(I_{n,a}=1)\le2/(a+1)\).  Let \(f\ge0\) be continuous with compact
support in \(J\).  Choose a step function
\(g=\sum_{\ell=1}^m c_\ell\one_{J_\ell}\), where the \(J_\ell\) are
disjoint half-open intervals, such that \(\|f-g\|_\infty\le\varepsilon\).
Since \(x\mapsto e^{-x}\) is one-Lipschitz on \([0,\infty)\),
\[
\left|
\Ee e^{-\langle f,\Xi_n\rangle}
-\Ee e^{-\langle g,\Xi_n\rangle}
\right|
\le\varepsilon\Ee\Xi_n(J)
\le C_J\varepsilon.
\tag{7.31}
\]
The already proved count-vector convergence gives convergence of the Laplace
functional for \(g\).  The same approximation applies to the limiting Poisson
process, whose expected mass on \(J\) is finite.  Letting first \(n\to\infty\)
and then \(\varepsilon\downarrow0\) proves convergence of Laplace functionals
for every nonnegative compactly supported continuous \(f\).  These Laplace
functionals determine the law in the vague topology, completing the proof.
\end{proof}

\begin{theorem}[Sukhatme fixed factorial moments]\label{thm:suk-fact}
For every fixed \(r\ge1\),
\[
        \Ee(F_n)_r\sim (e^{-1}\log n)^r.
\]
Consequently,
\[
        \frac{F_n}{e^{-1}\log n}\to1
\]
in probability.
\end{theorem}

\begin{proof}
Let
\[
        L_n=\lceil\log n\rceil,
        \qquad
        U_n=\lfloor n/\log n\rfloor,
\]
and put
\[
        G_n=\sum_{L_n\le a\le U_n}I_{n,a}.
\]
We first prove the factorial asymptotic for \(G_n\).  Fix \(\eta>0\).  In the expansion of \(\Ee(G_n)_r\), restrict to ordered distinct \(r\)-tuples whose logarithms are pairwise separated by at least \(\eta\).  For each such ordered tuple, sort the rates increasingly before applying Proposition \ref{prop:suk-joint}.  Pairwise logarithmic separation implies
\[
        \frac{a_{j+1}}{a_j}\ge e^\eta=1+(e^\eta-1),
\]
so Proposition \ref{prop:suk-joint} applies uniformly with separation parameter \(e^\eta-1\).  Hence the separated contribution is
\[
        e^{-r}
        \sum_{\substack{L_n\le a_1,\ldots,a_r\le U_n\\
                         a_i\ne a_j,\ |\log a_i-\log a_j|>\eta}}
        \frac1{a_1\cdots a_r}(1+o(1)),
\]
with the error uniform for fixed \(r\) and \(\eta\).  For fixed \(\eta\), the excluded logarithmic diagonal neighborhood has weighted lattice sum \(O_{r,\eta}((\log n)^{r-1})\).  Hence the separated weighted sum is
\[
        \left(\sum_{L_n\le a\le U_n}\frac1a\right)^r+o((\log n)^r).
\]
The non-separated part itself is negligible as well.  By Lemma \ref{lem:death-bound}, its contribution is bounded by a constant times the same logarithmic diagonal weighted sum, hence it is \(O_{r,\eta}((\log n)^{r-1})=o((\log n)^r)\).  Since
\[
        \sum_{L_n\le a\le U_n}\frac1a\sim \log n,
\]
we obtain
\[
        \Ee(G_n)_r\sim (e^{-1}\log n)^r.
\]

It remains to show that the omitted ranges do not affect the leading term.  By Lemma \ref{lem:death-bound}, any factorial-moment term containing at least one index outside \([L_n,U_n]\) is bounded by a constant times the corresponding product \(\prod 1/a_j\).  The logarithmic mass of the omitted set is
\[
        \sum_{a<L_n}\frac1a+
        \sum_{a>U_n}\frac1a
        =O(\log\log n)=o(\log n).
\]
Therefore the total contribution of tuples with at least one omitted index is
\[
        O((\log n)^{r-1}\log\log n)=o((\log n)^r).
\]
This proves \(\Ee(F_n)_r\sim(e^{-1}\log n)^r\).

For \(r=1,2\), the factorial moment asymptotics imply \(\Ee F_n\sim e^{-1}\log n\) and
\[
        \Ee F_n^2=\Ee(F_n)_2+\Ee F_n
        =(e^{-1}\log n)^2+o((\log n)^2).
\]
Consequently
\[
        \Var(F_n)=\Ee F_n^2-(\Ee F_n)^2=o((\log n)^2).
\]
Let \(\lambda_n=e^{-1}\log n\).  Since \(\Ee F_n/\lambda_n\to1\) and
\(\Var(F_n)=o(\lambda_n^2)\), Chebyshev's inequality gives
\(F_n/\lambda_n\to1\) in probability.
\end{proof}

% END ARCHIVED BLOCK: sukhatme-fixed-point-theory

\begin{thebibliography}{99}

\bibitem{Petrov1975}
V. V. Petrov.
\newblock \emph{Sums of Independent Random Variables}.
\newblock Ergebnisse der Mathematik und ihrer Grenzgebiete, Band 82.
\newblock Springer-Verlag, Berlin--Heidelberg, 1975.

\bibitem{Sukhatme1937}
P. V. Sukhatme.
\newblock Tests of significance for samples of the $\chi^2$-population with
two degrees of freedom.
\newblock \emph{Annals of Eugenics} 8 (1937), no. 1, 52--56.
\newblock \href{https://doi.org/10.1111/j.1469-1809.1937.tb02159.x}
{doi:10.1111/j.1469-1809.1937.tb02159.x}.

\bibitem{Renyi1953}
A. R\'enyi.
\newblock On the theory of order statistics.
\newblock \emph{Acta Mathematica Academiae Scientiarum Hungaricae} 4 (1953),
191--231.
\newblock \href{https://doi.org/10.1007/BF02127580}
{doi:10.1007/BF02127580}.

\end{thebibliography}

\end{document}
